\documentclass[11pt,a4paper]{amsart}

\usepackage[margin=1in]{geometry}
\usepackage{amsmath,amssymb,amsthm,mathtools}
\usepackage[expansion=false]{microtype}
\usepackage{enumitem}
\usepackage{graphicx}
\usepackage[colorlinks=true,linkcolor=black,citecolor=black,urlcolor=black]{hyperref}
\hypersetup{
  pdftitle={Boundary reconstruction and cofactor rigidity for skew Schur functions},
  pdfauthor={John Shields},
  pdfkeywords={skew Schur function, skew-equivalence, skew diagram reconstruction,
    Jacobi--Trudi determinant, Monge matrix, convex valuation,
    McNamara--van Willigenburg conjecture}
}
\pdfinfo{ /ORCID (https://orcid.org/0009-0004-3882-4571) }

\numberwithin{equation}{section}

\theoremstyle{plain}
\newtheorem{theorem}{Theorem}[section]
\newtheorem{proposition}[theorem]{Proposition}
\newtheorem{lemma}[theorem]{Lemma}
\newtheorem{corollary}[theorem]{Corollary}

\theoremstyle{definition}
\newtheorem{definition}[theorem]{Definition}
\newtheorem{example}[theorem]{Example}
\newtheorem{remark}[theorem]{Remark}

\DeclareMathOperator{\supp}{supp}
\DeclareMathOperator{\argmin}{argmin}

\newcommand{\Z}{\mathbb Z}
\newcommand{\R}{\mathbb R}
\newcommand{\Lam}{\Lambda}
\newcommand{\wt}{\nu}
\newcommand{\wto}{\nu_{\omega}}
\newcommand{\inin}{\operatorname{in}}

\title[Boundary reconstruction and cofactor rigidity]
{Boundary reconstruction and cofactor rigidity\\
for skew Schur functions}

\author{John Shields}
\date{September 11, 2026}
\thanks{Independent researcher, Ginza 2-15-2, KR Ginza II 5F, Chuo-ku, Tokyo, Japan 104-0061.\\
\hspace*{\parindent}Corresponding author: \href{mailto:math@johnnyshields.com}{math@johnnyshields.com}.  ORCID: \href{https://orcid.org/0009-0004-3882-4571}{0009-0004-3882-4571}.\\
\hspace*{\parindent}This work is archived at \href{https://doi.org/10.5281/zenodo.22723463}{doi:10.5281/zenodo.22723463}.}

\subjclass[2020]{Primary 05E05; Secondary 15A15, 90C24, 90C27}
\keywords{skew Schur function, skew-equivalence, skew diagram reconstruction, Jacobi--Trudi determinant, Monge matrix, convex valuation, McNamara--van Willigenburg conjecture}

\begin{document}

\begin{abstract}
Two boundary observations determine a connected skew diagram up to translation, with neither boundary prescribed: its skew Schur function together with that of its first-row deletion, or of a northeast row extension of any positive length.  Two extensions of distinct positive lengths also suffice, with neither the lengths nor the diagram's own function supplied.  Strict convexity makes the Jacobi--Trudi staircase strictly Monge, so every minor an increasing matching can fill has a unique leading term; successive prefix minors eliminate cofactor coordinates, and the quotients of their leading monomials return the row and column parameters one row at a time.  For equal diagrams with a common inner boundary the cofactor difference has a rigid adjacent-edge fingerprint, whose supported prefix has single-monomial initial forms that do not move with the valuation, for every admissible strictly convex weight and with no bound on the number of changed rows or supported shifts.  That chain has the local adjacency pattern of McNamara--van Willigenburg composition, with factor reversal the orientation freedom that pattern would leave; no factorization is claimed.
\end{abstract}

\maketitle

\section{Introduction}

McNamara and van Willigenburg predict that skew Schur equality is generated by a composition operation together with reversal of its factors \cite{McNamara2009Classification}.  On ribbons this is a theorem \cite{Billera2006DecomposableCompositions}, and the degenerate case in which one side is an ordinary Schur function is settled \cite{VanWilligenburg2005Equality}, as is the multiplicity-free case \cite[Thm.~3.14]{Gutschwager2009MultiplicityFree}; the general necessity problem remains open, and \cite{McNamara2014ComparingSkew} surveys what is known.  We reconstruct both boundaries of a connected diagram from two skew Schur functions, and we study equality with the inner boundary prescribed, imposing no bound on the number of outer-boundary positions that change.  The local single-junction situation, in which the two outer boundaries differ at exactly one junction,
\[
\lambda'_p=\lambda_p+1,
\qquad
\lambda'_{p+1}=\lambda_{p+1}-1,
\]
lies in that class but is not what the theorems require.  What we prove is cofactor rigidity and boundary determination, not a classification of equalities by composition factors; nor do the results subsume the distinct transverse class treated by Jin and Li \cite{Jin2024ThickenedRibbon}.  A different enrichment is studied by Jin and van Willigenburg \cite[Thm.~3.1]{Jin2025NoncommutingEquality}, who classify equality of skew Schur functions in noncommuting variables; those functions map onto the ordinary ones \cite[Lem.~2.19]{Jin2025NoncommutingEquality}, so equality of the richer functions is a stronger condition and not a classification of the ordinary one.


\subsection{Main results}\label{subsec:intro-main}

Our first result determines a diagram from two boundary observations.
Write \(\bar X\) for the diagram obtained by deleting the entire first row
of \(X\), and \(X\cdot(a)\) for its northeast extension by a row of length
\(a\) meeting the old top row in exactly one column.  These are two
different boundary operations (Figure~\ref{fig:boundary-observations}).

\begin{theorem}[Reconstruction from two boundary observations]\label{thm:boundary-reconstruction}
On translation classes of connected skew diagrams, each of the maps
\[
X\longmapsto(s_X,s_{\bar X}),
\qquad
X\longmapsto(s_X,s_{X\cdot(a)})
\quad(a\ge1\text{ fixed})
\]
is injective.  Here \(s_\varnothing=1\).
\end{theorem}

The second map needs \(a\) only to be well defined, not supplied: \(s_X\) is homogeneous of degree \(|X|\), so \(\deg s_{X\cdot(a)}-\deg s_X=a\), and the unordered pair \(\{s_X,s_{X\cdot(a)}\}\) already says which observation is which and what \(a\) is.  Collisions between different lengths are excluded by that degree count alone.

\begin{corollary}[Reconstruction from two positive extensions]\label{cor:two-positive-extensions}
Let \(a,b>0\) be distinct.  On translation classes of connected skew diagrams, the unordered pair
\[
\bigl\{s_{X\cdot(a)},s_{X\cdot(b)}\bigr\}
\]
determines \(X\) and determines \(\{a,b\}\).  The lengths need not be supplied.
\end{corollary}

Two distinct positive extensions therefore suffice without observing
\(s_X\) itself.  Both statements are proved in Section~\ref{sec:reconstruction}, and the
extension lengths are arbitrary subject to the conditions stated above.

\begin{figure}[htbp]
\centering
\includegraphics[width=\linewidth]{figures/monge-boundary-observations.pdf}
\caption{The two boundary observations of Theorem~\ref{thm:boundary-reconstruction}
on the equal pair of Example~\ref{ex:actual-equal-chain}:
\(X=(4,4,4,2)/\mu\) and \(Y=(4,4,3,3)/\mu\) with \(\mu=(3,2,1,0)\).
Dotted cells belong to \(\mu\) and not to the diagram.  Every panel is drawn
in one grid, so a row that is deleted or attached is visible as such.  Hatching
marks the first row in the left column and the attached row in the right.  The
attached row meets the old top row in exactly one column, which is the contact
the definition of \(X\cdot(a)\) requires.  Although \(s_X=s_Y\), each
observation separates the pair.}
\label{fig:boundary-observations}
\end{figure}

Theorem~\ref{thm:boundary-reconstruction} is proved by recovering successive
row and column parameters from two cofactors.
Write \(C(X)\) for the signed maximal-cofactor vector of the capped
Jacobi--Trudi matrix, defined in Section~\ref{sec:coordinates}.  For distinct
connected \(X=\lambda/\mu\) and \(Y=\kappa/\mu\) with \(n\) nonempty rows
and \(s_X=s_Y\), equality itself aligns the cap grids: the largest generator
index occurring in \(s_X\) is \(\lambda_0-\mu_{n-1}+n-1\) by
\eqref{eq:largest-index}, so \(s_X\) determines \(\lambda_0\).  From \(s_X\) alone, \cite[Cor.~8.11]{Reiner2007Coincidences} recovers the row
overlap partitions, and \cite[Thm.~8.10]{Reiner2007Coincidences} recovers the
skew Schur function of the diagram obtained by removing the top cell of every
column, which is not the first-row deletion.  The ordered row lengths and adjacent overlaps do determine a connected
diagram up to translation, but they cannot in general be recovered from its
skew Schur function alone \cite[discussion preceding Lem.~8.5]{Reiner2007Coincidences};
this is why a second observation is needed at all.  What the cofactor
observations add is exactly those ordered parameters, and with them both
boundaries.  Even the
overlap hierarchies of both observed diagrams do not suffice: \(U=(7,7,5,4)/(6,4,2,0)\)
and \(V=(7,7,6,3)/(6,4,2,0)\) satisfy \(s_U=s_V\) and agree in every row
overlap partition, before and after first-row deletion, yet
\(s_{\bar U}-s_{\bar V}=h_3h_7-h_4h_6\ne0\).

For equal pairs with a common inner boundary the cofactor difference satisfies
the further rigidity stated in Theorem~\ref{thm:equal-pair-rigidity}, at the
point where it is proved.  Using zero-based
row and column indices, if \(v\) is the first changed row and
\(G=C(X)-C(Y)\), then \(G_0=0\), while in each common-row relation
\(0\le i<v\) the minimum valuation occurs at exactly the adjacent columns
\(i+1\) and \(i+2\).  This holds for every admissible strictly convex
weight.  Moreover, the initial forms of \(G_1,\ldots,G_{v+1}\) are single
monomials with unit coefficients independent of the weight
(Proposition~\ref{prop:monomial-prefix}).  In particular \(G_1\ne0\), and
this first positive coordinate is what the boundary observations detect.

A second branch concerns the ratio identity \eqref{eq:ratio-identity}, a
condition on capped functions rather than on equality of the parents.
For connected parents with a common inner boundary, it supplies equal
row-length multisets, a virtual row and the full adjacent-edge chain
(Theorem~\ref{thm:ratio-rigidity}); two ratio identities at the same base force the
defect to vanish (Corollary~\ref{cor:double-cap-kills}), and a nonzero defect
admits at most one ratio identity in the entire parameter plane
(Corollary~\ref{cor:unique-ratio}).  Given the full cofactor vector rather than two
observations, the diagram can be read off coarser data still: two scalars per
column, the ordinary degree and the square valuation of each cofactor
(Proposition~\ref{prop:stable-cofactor}).  The two branches are distinct:
no distinct equal pair with a common inner boundary satisfies a ratio
identity (Corollary~\ref{cor:boundary-separation}).

\subsection{Proof strategy}

Equal polynomials have the same minimum under every additive monomial
valuation, so valuating \(s_X=s_Y\) itself cannot distinguish the diagrams.
Nor does the square-weight collapse faithfully record every nonzero defect
coordinate; Section~\ref{sec:scope} gives a concrete example.  We instead
valuate the nonzero coordinates of the cofactor defect and use exact
kernel relations to constrain where their minimum can occur, rather than to
single out the coefficients that survive it.

The capped matrix \(M_{ij}=h_{d_j-A_i}\),
defined in Section~\ref{sec:coordinates}, has increasing row and column
parameters.  Strict convexity makes crossed assignments more expensive than
monotone ones, so every feasible minor has a unique monotone leading
matching.  For square weight the normalized cofactor actions
\(\sigma_j^X=\wt(C_j(X))+d_j^2\) satisfy an exact recurrence whose increment is
\(2A_j^X(d_{j+1}-d_j)\).

With neither boundary prescribed, that uniqueness already reconstructs.
Prefix minors eliminate the cofactor coordinates one at a time: each residual
has a unique initial monomial, and the quotient of two consecutive ones
returns the next row parameter together with the next cap shift, the two
numerator indices straddling the denominator strictly because the row
parameters increase, the cap grid increases, and the diagram is connected.
That is Theorem~\ref{thm:boundary-reconstruction}, and it uses neither the
ratio identity nor a comparison between two diagrams.

For two diagrams sharing an inner boundary, equal row-length data synchronizes
the two parent profiles through their common prefix, and the first changed row
\(v\) produces a strict split at \(q=v+1\).
Exact common-row relations require tied minima.  Strict-Monge single
crossing makes their rightmost locations march to the right, while the
first-split anchor prevents columns beyond \(q\) from minimizing.  For an
actual equal pair, \(G_0=s_X-s_Y=0\) removes column zero from the minimum.
The available columns then force the chain
\begin{equation}\label{eq:edge-chain-intro}
\{1,2\},\ \{2,3\},\ldots,\ \{v,v+1\}.
\end{equation}
No ratio identity is used in this branch either.  Its first positive
coordinate cannot vanish, giving first-row-deletion separation; every
positive cap isolates that same coordinate at its lowest weight, giving
extension separation.

On the ratio-identity locus, the ratio identity supplies both cap
normalization and one virtual row before the genuine top row.  Its extra
minimizing pair \(\{0,1\}\) precedes the chain
\eqref{eq:edge-chain-intro}.  The square-weight consequence is
\emph{Cancellative Leftmost}: once the leading cancellation at the top row
is complete, the leftmost positive supported shift is one of the minimizers.
The iterated-cap identity \eqref{eq:iterated-cap} then gives
Corollary~\ref{cor:two-positive-extensions} directly from boundary
reconstruction.

\subsection{Organization}

Section~\ref{sec:coordinates} fixes the cofactor coordinates, the exact
relations, and the class of admissible convex weights every later hypothesis
names.  Section~\ref{sec:geometry} establishes the Strict-Monge minor
geometry and action recurrence.  Section~\ref{sec:reconstruction} proves the
two reconstruction results stated above, and recovers a diagram from two
scalars per cofactor.  Section~\ref{sec:localization} localizes
the defect to the first split, proves cap-fixed adjacent-edge rigidity
there, and applies it to actual equality.
Section~\ref{sec:full-rigidity} drops the common cap grid: the ratio
identity supplies that grid and a virtual row, which together give full
adjacent-edge rigidity.  Section~\ref{sec:two-cap} proves two-cap
exclusion and uniqueness of the ratio identity.
Section~\ref{sec:application} shows that every positive cap separates an
equal pair, and reads the cancellation chain against composition and
reversal.  Section~\ref{sec:scope} records the limitations and the
missing steps toward classification.  Appendix~\ref{app:structural} collects
independent structural consequences, including the abstract Strict-Monge
principle, and Appendix~\ref{app:rank4} treats the passive tail.

\section{Capped cofactors and admissible convex weights}\label{sec:coordinates}

The geometric argument begins in
Section~\ref{sec:geometry}; its equality and ratio-identity consequences are
stated alongside their proofs in Sections~\ref{sec:localization} and
\ref{sec:full-rigidity}.  The ratio identity is introduced here for later
use, but nothing before Section~\ref{sec:full-rigidity} assumes it.

\subsection{Conventions}

We work in \(\Lam=\Z[h_1,h_2,\ldots]\),
with
\[
h_0=1,
\qquad
h_k=0\quad(k<0).
\]
The generators \(h_1,h_2,\ldots\) are algebraically independent.

We use zero-based row indices throughout this paper.  Let \(X=\lambda/\mu\)
be a connected skew diagram with \(n\) nonempty rows indexed by \(i=0,\ldots,n-1\).
Pad \(\mu\) by zeros to length \(n\), and put \(L=\lambda_0\).
Define
\[
a_i=\lambda_i-(i+1),
\qquad
b_0=L-1,
\qquad
b_j=\mu_{j-1}-j\quad(1\le j\le n).
\]
Set
\begin{equation}\label{eq:Adef}
A_i=b_0-a_i=L-\lambda_i+i,
\end{equation}
and
\begin{equation}\label{eq:ddef}
d_j=b_0-b_j.
\end{equation}
Thus
\begin{equation}\label{eq:dexplicit}
d_0=0,
\qquad
 d_j=L-\mu_{j-1}+j-1\quad(j\ge1).
\end{equation}
The tuple \((d_0,\ldots,d_n)\) is the \emph{cap grid} of \(X\); by
\eqref{eq:dexplicit} it is determined by \(\mu\) and \(L\) alone.
The capped cofactor matrix is
\begin{equation}\label{eq:capmatrix}
M_{ij}=h_{d_j-A_i}.
\end{equation}
Writing a Jacobi--Trudi entry as a difference of cumulative sums of row lengths and overlaps, as \eqref{eq:Adef} and \eqref{eq:dexplicit} do, is the presentation of \cite[(1.11)--(1.12)]{King2008SkewPositivity}; the row-overlap coordinates themselves are those of \cite[\S8.2]{Reiner2007Coincidences}.

The signed maximal cofactor obtained by deleting column \(j\) is
\begin{equation}\label{eq:cofactor}
C_j(X)=(-1)^j\det M_{\widehat j}.
\end{equation}
When it is useful to label a coordinate by its cap shift rather than its column, we write \(C_{d_j}(X)\), and we extend it by zero off the grid: \(C_d(X)=C_j(X)\) if \(d=d_j^X\) for some \(j\), and \(C_d(X)=0\) otherwise.

For two parents \(A,B\) define
\begin{equation}\label{eq:defect}
G_d=C_d(A)-C_d(B),
\end{equation}
which is defined on the union of the two cap grids and needs no common grid; on a common grid it is the column-indexed \(G_j=C_j(A)-C_j(B)\), and the two notations are used interchangeably there.  Extension by zero does not change \(\supp G\), since a vanishing coordinate is not in the support.
The shifted cofactor combination is
\begin{equation}\label{eq:Pk}
P_k(G)=\sum_d G_d h_{k+d},
\end{equation}
the sum running over the union of the two grids; on a common grid it is \(\sum_j G_j h_{k+d_j}\).

Let \(v\ge1\) be the first changed row:
\begin{equation}\label{eq:first-change}
A_i^A=A_i^B\quad(i<v),
\qquad
A_v^A\ne A_v^B.
\end{equation}
Put
\begin{equation}\label{eq:q-def}
q=v+1.
\end{equation}

\subsection{Exact cap relations}

For any \(n\times(n+1)\) matrix, the signed vector of maximal cofactors lies in its right kernel.  Hence, for every row common to the two parents,
\begin{equation}\label{eq:common-row-relation}
\sum_j h_{d_j-A_i}G_j=0.
\end{equation}
In particular this holds for all \(0\le i<v\).

The same kernel relation applied to a single parent at its own row \(i\) reads \(P_{-A_i}(C(X))=0\),
since the cap entry at shift \(-A_i\) coincides entrywise with row \(i\).  Because \(A_0=0\) by \eqref{eq:Adef}, the case \(i=0\) is
\begin{equation}\label{eq:P0-vanishes}
P_0(C(X))=0,
\end{equation}
which holds separately for each connected parent.  Thus \(P_0(G)=0\) without assuming that the parents are skew-equivalent; both roles of the ratio identity use this same relation.

The ratio identity at \((r,a)\) is the relation
\begin{equation}\label{eq:ratio-identity}
P_{r+a}(G)=h_aP_r(G),
\qquad r\ge0,\quad a>0.
\end{equation}

Although \eqref{eq:ratio-identity} is written in cofactor coordinates, it is a statement about the two diagrams.  Write \(X\cdot(k)\) for \(X\) with a row of length \(k\) attached to the northeast, meeting its top row in exactly one column.  Expanding the determinant of \(X\cdot(k)\) along that row gives
\begin{equation}\label{eq:Pk-is-capped}
P_k(C(X))=s_{X\cdot(k)}
\qquad(k\ge1),
\end{equation}
so, for \(r\ge1\), \eqref{eq:ratio-identity} reads
\[
h_as_{A\cdot(r)}-s_{A\cdot(r+a)}
=
h_as_{B\cdot(r)}-s_{B\cdot(r+a)} :
\]
the two capped parents have equal signed Pieri remainders.  At \(r=0\) the attachment \(X\cdot(0)\) is not defined, and \eqref{eq:ratio-identity} reads instead \(P_a(G)=0\), that is \(s_{A\cdot(a)}=s_{B\cdot(a)}\), since \eqref{eq:P0-vanishes} gives \(P_0(G)=0\).  The standing hypothesis of Section~\ref{sec:full-rigidity} is therefore a skew Schur condition on the original objects rather than a relaxation of one.  The signed remainder is itself the skew Schur function of an iterated extension.

\begin{lemma}\label{lem:extension-cofactors}
Let \(X\) be connected and let \(k>0\).  The cap grid of \(X\cdot(k)\) is \((0,k,k+d_1,\ldots,k+d_n)\), and in column order
\[
C(X\cdot(k))
=
\bigl(P_k(C(X)),-C_0(X),\ldots,-C_n(X)\bigr).
\]
Consequently, for \(r,a>0\),
\begin{equation}\label{eq:iterated-cap}
s_{(X\cdot(r))\cdot(a)}
=
h_as_{X\cdot(r)}-s_{X\cdot(r+a)}.
\end{equation}
\end{lemma}

\begin{proof}
The row parameters of \(X\cdot(k)\) are \((0,k+A_0,\ldots,k+A_{n-1})\), so its capped matrix is
\[
\begin{pmatrix}
1&(h_{k+d_j})_{j=0}^{n}\\
0&M(X)
\end{pmatrix}.
\]
Deleting column zero leaves the lower block bordered by its top row, which is \(P_k(C(X))\) by the definition \eqref{eq:Pk}.  Deleting column \(j+1\) and expanding along the first column leaves \(-C_j(X)\), the sign being the one the cofactor carries.  Applying \(P_a\) to the resulting vector at the shifted grid gives
\[
P_a(C(X\cdot(r)))
=
h_aP_r(C(X))-\sum_j C_j(X)h_{a+r+d_j},
\]
which is \eqref{eq:iterated-cap} by \eqref{eq:Pk-is-capped}.
\end{proof}

So for \(r>0\) the ratio identity is exactly \(s_{(A\cdot(r))\cdot(a)}=s_{(B\cdot(r))\cdot(a)}\), and at \(r=0\) it is \(s_{A\cdot(a)}=s_{B\cdot(a)}\).  These are equalities of extended diagrams, not assertions about the parents.  The scalar identity \eqref{eq:iterated-cap} is the classical product rule for concatenation and near-concatenation \cite[Ch.~I, \S5, Ex.~21(a)]{Macdonald1995SymmetricFunctions}, stated for skew diagrams in \cite[Prop.~4.1]{Reiner2007Coincidences}; what Lemma~\ref{lem:extension-cofactors} adds is the cofactor transform it comes from.  Either way it is a determinant identity and makes no claim about multisets of Pieri shapes.

\subsection{Weights and tropical notation}

\begin{definition}\label{def:convex-weight}
An \emph{admissible convex weight} is a function
\[
\omega:\Z_{\ge0}\to\R
\]
such that
\[
\omega(0)=0,
\]
the increments
\[
\Delta\omega(k)=\omega(k)-\omega(k-1)
\]
are strictly increasing for \(k\ge1\), and
\[
\Delta\omega(k)>0
\qquad(k\ge1).
\]
\end{definition}

The positivity of the increments is unnecessary for the cap-fixed Strict-Monge theorem itself, but it is convenient when the ratio identity is converted into a virtual row in Subsection~\ref{subsec:virtual}.  The principal example is \(\omega(k)=k^2\).  Throughout, \(\omega\) denotes such a weight and never the involution \(s_{\lambda/\mu}\mapsto s_{(\lambda/\mu)^{t}}\) on \(\Lam\), nor a ribbon.

For an \(h\)-monomial
\[
m=\prod_{r=1}^{\ell}h_{\eta_r},
\]
define
\[
\wto(m)=\sum_{r=1}^{\ell}\omega(\eta_r).
\]
For a nonzero polynomial \(f=\sum_m c_m m\), define \(\wto(f)=\min_{c_m\ne0}\wto(m)\), and put \(\wto(0)=+\infty\).

Define
\begin{equation}\label{eq:alpha-general}
\alpha_j^{\omega}
=
\wto(G_j)+\omega(d_j),
\end{equation}
with \(\alpha_j^\omega=+\infty\) when \(G_j=0\).

For a common-prefix row \(i<v\), define
\begin{equation}\label{eq:F-general}
F_i^\omega(j)
=
\alpha_j^\omega
+
\omega(d_j-A_i)
-
\omega(d_j),
\end{equation}
whenever \(G_j\ne0\) and \(d_j\ge A_i\), and set \(F_i^\omega(j)=+\infty\) otherwise.  This is exactly the valuation of the \(j\)-th nonzero summand in \eqref{eq:common-row-relation}.

Let \(\mathcal M_i^\omega=\argmin_j F_i^\omega(j)\).

Define \(m_+^\omega=\min_{j>0}\alpha_j^\omega\)
and \(T_{0,\omega}^+=\{d_j:j>0,\ \alpha_j^\omega=m_+^\omega\}\).

The ratio identity also contributes one row parameter to the left of the genuine top row, written \(-\gamma\) and fixed by \eqref{eq:gamma}; \(\mathcal M_{-1}^{\omega}\) denotes the minimizer set of that virtual row.  Section~\ref{subsec:virtual} derives both.

\section{Strict-Monge geometry of cofactor actions}\label{sec:geometry}

\subsection{Ordered grids and connectedness}\label{subsec:setup}

\begin{lemma}\label{lem:grid-order}
For a skew diagram as above,
\[
0=A_0<A_1<\cdots<A_{n-1}
\]
and
\[
0=d_0<d_1<\cdots<d_n.
\]
\end{lemma}

\begin{proof}
From \eqref{eq:Adef}, \(A_{i+1}-A_i =\lambda_i-\lambda_{i+1}+1>0\).
Similarly, \(d_1=L-\mu_0>0\), and for \(j\ge1\), \(d_{j+1}-d_j=\mu_{j-1}-\mu_j+1>0\).
\end{proof}

\begin{lemma}\label{lem:diagonal-feasible}
If \(X\) is connected, then
\begin{equation}\label{eq:diagonal-feasible}
A_i\le d_i
\qquad(0\le i<n).
\end{equation}
\end{lemma}

\begin{proof}
For \(i=0\) both sides vanish.  For \(i\ge1\), \(d_i-A_i=\lambda_i-\mu_{i-1}-1\).
Connectedness of rows \(i-1\) and \(i\) gives \(\lambda_i>\mu_{i-1}\), hence the right-hand side is nonnegative.
\end{proof}

\begin{remark}
Topological connectedness is not used directly; it is used through the matrix inequality \eqref{eq:diagonal-feasible}, which guarantees feasibility of monotone determinant matchings.
\end{remark}

\subsection{Cofactors as skew Schur functions}

\begin{lemma}\label{lem:cofactor-factorization}
Let \(X=\lambda/\mu\) have \(n\) nonempty rows, not necessarily connected.  For \(0\le j\le n\) put
\begin{equation}\label{eq:nu-j}
\nu^{(j)}
=
\bigl(\lambda_0,\ \mu_0+1,\ \ldots,\ \mu_{j-2}+1,\ \mu_j,\ \mu_{j+1},\ \ldots,\ \mu_{n-1}\bigr)
\qquad(j\ge1),
\qquad
\nu^{(0)}=\mu.
\end{equation}
Then
\begin{equation}\label{eq:cofactor-support}
C_j(X)\ne0
\iff
\lambda_i>\mu_{i-1}
\quad\text{for every }1\le i<j,
\end{equation}
and where the right-hand side holds, \(\nu^{(j)}\) is a partition contained in \(\lambda\) and
\begin{equation}\label{eq:cofactor-schur}
C_j(X)=(-1)^j s_{\lambda/\nu^{(j)}}.
\end{equation}
The condition in \eqref{eq:cofactor-support} is empty for \(j\le1\) and is exactly connectedness at \(j=n\), so
\[
C_n(X)\ne0
\iff
X\text{ is connected},
\]
while for connected \(X\) it holds at every \(j\) and no cofactor vanishes.
\end{lemma}

\begin{proof}
Delete column \(j\) from \(M\) and index the surviving columns by \(t=0,\ldots,n-1\), so that column \(t\) carries the parameter \(b_t\) for \(t<j\) and \(b_{t+1}\) for \(t\ge j\).  Its entry in row \(i\) is \(h_{a_i-b_{\bullet}}\), and \(a_i=\lambda_i-i-1\), so the surviving matrix is the Jacobi--Trudi matrix \((h_{\lambda_i-\nu_t-i+t})\) of \(\lambda/\nu\) exactly when \(\nu_t\) equals the surviving parameter plus \(t+1\).  Reading that off gives \(\nu_0=b_0+1=\lambda_0\) for \(j\ge1\), then \(\nu_t=\mu_{t-1}+1\) for \(1\le t<j\), and \(\nu_t=\mu_t\) for \(t\ge j\), which is \eqref{eq:nu-j}; for \(j=0\) every column satisfies \(t\ge j\) and \(\nu^{(0)}=\mu\).

The entries of \(\nu^{(j)}\) are weakly decreasing because \(\mu\) is and because \(\mu_{j-2}+1\ge\mu_j\).  Containment \(\nu_t\le\lambda_t\) is \(\lambda_0\le\lambda_0\) at \(t=0\), is \(\mu_t\le\lambda_t\) for \(t\ge j\), and for \(1\le t<j\) is \(\mu_{t-1}+1\le\lambda_t\) --- which is the right-hand side of \eqref{eq:cofactor-support}, and the only place the hypothesis is consumed.  Hence \(\lambda/\nu^{(j)}\) is a skew diagram, \(\det M_{\widehat\jmath}=s_{\lambda/\nu^{(j)}}\) is nonzero, being positive in the monomial symmetric function basis, and \eqref{eq:cofactor-schur} follows from \eqref{eq:cofactor}.

For the converse, suppose \(\lambda_i\le\mu_{i-1}\) for some \(1\le i<j\).  Then \eqref{eq:Adef} and \eqref{eq:dexplicit} give \(d_i-A_i=\lambda_i-\mu_{i-1}-1<0\), so \(A_i>d_i\).  The \(i+1\) columns \(0,\ldots,i\) are all retained, since \(i<j\), and for \(t\le i\) and \(k\ge i\) we have \(d_t\le d_i<A_i\le A_k\) by Lemma~\ref{lem:grid-order}, so every one of their entries vanishes in every row \(i,\ldots,n-1\).  Those \(i+1\) columns are therefore supported in the \(i\) rows \(0,\ldots,i-1\) alone, Hall's condition fails, and every term of \(\det M_{\widehat\jmath}\) contains a zero entry.
\end{proof}

Connectedness therefore enters twice in the same way: through \eqref{eq:diagonal-feasible} and through the containment \(\nu^{(j)}\subseteq\lambda\), and in both cases only as \(\lambda_i>\mu_{i-1}\).  That a minor of a banded Toeplitz matrix is a skew Schur polynomial is the dictionary of \cite{Maximenko2017TrenchFormulas}; \eqref{eq:nu-j} is the instance of it that names the shape belonging to a maximal cap cofactor.

\begin{remark}
Equality of row-length multisets is an invariant of skew Schur equality: it is the \(k=1\) case of the row-overlap hierarchy of \cite[Cor.~8.11]{Reiner2007Coincidences}, restated as \cite[Thm.~3.5]{McNamara2014ComparingSkew}.
\end{remark}

\subsection{Convex valuations}\label{subsec:valuations}

\begin{lemma}\label{lem:valuation-laws}
For nonzero \(f,g\in\Lam\),
\[
\wto(fg)=\wto(f)+\wto(g).
\]
For arbitrary \(f,g\),
\[
\wto(f+g)\ge\min\{\wto(f),\wto(g)\}.
\]
If the two minimum valuations differ, equality holds.
\end{lemma}

\begin{proof}
The weight defines a grading of the polynomial ring.  The product of the two nonzero initial forms is nonzero because \(\Lam\) is an integral domain.  The assertion for sums is immediate.
\end{proof}

\subsection{Strict-Monge uncrossing}

\begin{lemma}[Strict-Monge inequality]\label{lem:strict-monge}
Let \(A<A'\) and \(d<d'\), with all four differences in the expression below nonnegative.  Then
\[
\omega(d-A)+\omega(d'-A')
<
\omega(d'-A)+\omega(d-A').
\]
\end{lemma}

\begin{proof}
Write
\[
x=d-A',\qquad p=A'-A>0,\qquad \zeta=d'-d>0.
\]
Then the crossed cost minus the aligned cost is \(\omega(x+p+\zeta)+\omega(x)-\omega(x+p)-\omega(x+\zeta)\).
This can be written as
\[
\sum_{u=1}^{p}\bigl(\Delta\omega(x+\zeta+u)-\Delta\omega(x+u)\bigr),
\]
which is strictly positive because \(\Delta\omega\) is strictly increasing.
\end{proof}

\begin{lemma}\label{lem:feasible-uncrossing}
Suppose \(A<A'\) and \(d<d'\).  If the crossed lower-left entry \(h_{d-A'}\) is nonzero, then so are the aligned entries
\[
h_{d-A},\qquad h_{d'-A'}
\]
and the remaining crossed entry \(h_{d'-A}\).
\end{lemma}

\begin{proof}
The crossed lower-left entry being nonzero gives \(d\ge A'\).  Hence \(d\ge A\) and \(d'>d\ge A'\), proving feasibility of both aligned entries, and \(d'>d\ge A'>A\) gives the remaining crossed entry.
\end{proof}

\begin{theorem}[Strict-Monge minor theorem]\label{thm:monge-minor}
Consider any square submatrix of a staircase matrix
\[
(h_{d_j-A_i})
\]
whose selected row parameters and selected column parameters are in increasing order.  Call the matching \emph{feasible} when every entry it selects is nonzero, that is when each selected column parameter is at least its matched row parameter.  If the increasing row-column matching is feasible, then it is the unique minimum-\(\omega\)-weight determinant term.  In particular the minor is nonzero and its \(\omega\)-initial monomial has coefficient \(+1\), the increasing matching being the identity permutation of the selected index sets.  Conversely, if that matching is not feasible the minor vanishes, so the minor is nonzero precisely when it is feasible.
\end{theorem}

\begin{proof}
Any non-increasing surviving permutation has an inversion.  By Lemma~\ref{lem:feasible-uncrossing} the inversion can be uncrossed without destroying feasibility, and by Lemma~\ref{lem:strict-monge} the uncrossing strictly lowers total \(\omega\)-weight.  Repeated uncrossing reaches the increasing permutation and strictly lowers the weight at every step.

The converse is the same computation read backwards.  Lemma~\ref{lem:feasible-uncrossing} turns a surviving crossed pair into a surviving aligned one, so if any permutation survives then so does the increasing one, and the increasing matching is feasible.  A minor whose increasing matching is infeasible therefore has no surviving permutation at all, and every determinant term vanishes.
\end{proof}

Uncrossing an inversion against a Monge cost is classical assignment theory; \cite{Burkard1996MongeProperties} surveys it.  What is specific to the present setting is that the Jacobi--Trudi staircase supplies the Monge structure through a convex valuation of the generators rather than through numerical entries.

\begin{remark}
A Jacobi--Trudi determinant contains many signed permutation terms, but its lowest convex geometry is the unique monotone transport between two ordered sets.
\end{remark}

\subsection{Exact initial monomials}\label{subsec:action}

Delete cap column \(j\).  In the increasing matching, row \(i\) uses
\begin{equation}\label{eq:matching-index}
k_i(j)=
\begin{cases}
 i,&i<j,\\
 i+1,&i\ge j.
\end{cases}
\end{equation}
Since \(k_i(j)\ge i\), Lemma~\ref{lem:diagonal-feasible} makes this matching feasible.

\begin{proposition}\label{prop:cofactor-initial}
For every connected \(X\) and every cap column \(j\),
\[
C_j(X)\ne0,
\]
and
\begin{equation}\label{eq:cofactor-initial}
\inin_{\omega}C_j(X)
=
(-1)^j
\left(\prod_{i<j}h_{d_i-A_i^X}\right)
\left(\prod_{i\ge j}h_{d_{i+1}-A_i^X}\right).
\end{equation}
Consequently
\begin{equation}\label{eq:cofactor-val}
\wto(C_j(X))
=
\sum_{i<j}\omega(d_i-A_i^X)
+
\sum_{i\ge j}\omega(d_{i+1}-A_i^X).
\end{equation}
\end{proposition}

\begin{proof}
Apply Theorem~\ref{thm:monge-minor} to the maximal minor obtained by deleting column \(j\).
\end{proof}

Factors \(h_0=1\) cause no difficulty: they contribute zero to the weight and may simply be omitted from the visible monomial.

At \(j=0\) the monomial \eqref{eq:cofactor-initial} is the diagonal term of the full Jacobi--Trudi determinant, and that its coefficient is \(+1\) and that it is dominance-minimal is \cite[Prop.~6.2(ii)]{Reiner2007Coincidences}, proved there by the same uncrossing step.  That statement is the stronger one and it already covers every admissible convex weight: monomials of a homogeneous determinant have equal total degree, so a dominance-minimal one minimizes \(\sum_j\omega(e_j)\) for every convex \(\omega\) by the Karamata inequality \cite[Ch.~3, \S C]{MarshallOlkinArnold2011Majorization}.  Combined with Lemma~\ref{lem:cofactor-factorization}, which identifies each cofactor as a skew Schur function, it also gives \eqref{eq:cofactor-initial} for every maximal minor.  We give the direct Strict-Monge proof because it produces the action recurrence and the single-crossing comparisons that Sections~\ref{sec:localization} and \ref{sec:full-rigidity} run on, not because its conclusion for one minor is new.  What is new here concerns the \emph{difference} of two cofactors after the parent initial terms cancel: the exact row relations and the first-split anchor then pin both the minimizing locations and the surviving initial monomials.

\subsection{Normalized action}

Define
\begin{equation}\label{eq:sigma-def}
\sigma_j^{X,\omega}
=
\wto(C_j(X))+\omega(d_j).
\end{equation}

\begin{proposition}\label{prop:action-recurrence}
For \(0\le j<n\),
\[
\begin{aligned}
\sigma_{j+1}^{X,\omega}-\sigma_j^{X,\omega}
&=
\omega(d_j-A_j^X)-\omega(d_{j+1}-A_j^X)\\
&\quad+\omega(d_{j+1})-\omega(d_j).
\end{aligned}
\]
For fixed \(d_j<d_{j+1}\), the right-hand side is strictly increasing as a function of \(A_j^X\).
\end{proposition}

\begin{proof}
The increasing matchings for deletion of \(j\) and \(j+1\) agree in every row except row \(j\), yielding the displayed difference.  Strict increase in \(A_j^X\) is an immediate form of the Strict-Monge inequality.
\end{proof}

For the square weight \(\omega(k)=k^2\) the recurrence collapses to \(\sigma_{j+1}^X-\sigma_j^X =2A_j^X(d_{j+1}-d_j)\).
This is the discrete action formula used in the geometric interpretation.

\subsection{Initial row energy}

Let \(\rho_i^X=\lambda_i^X-\mu_i\)
be the row lengths, so that deletion of cap column \(0\) matches row \(i\) to \(d_{i+1}-A_i^X=\rho_i^X\).

\begin{proposition}\label{prop:initial-energy}
\begin{equation}\label{eq:initial-energy}
\sigma_0^{X,\omega}=\sum_i\omega(\rho_i^X).
\end{equation}
In particular equal row-length multisets imply
\[
\sigma_0^{A,\omega}=\sigma_0^{B,\omega}
\]
for every admissible convex weight \(\omega\).
\end{proposition}

\begin{proof}
By \eqref{eq:cofactor-val} at \(j=0\), \(\sigma_0^{X,\omega}=\sum_i\omega(d_{i+1}-A_i^X)\), and \(d_{i+1}-A_i^X=\rho_i^X\).  Equal row-length multisets therefore give equal sums.
\end{proof}

\begin{remark}
For one fixed weight \(\omega\), the argument only needs equality of the single convex energy in \eqref{eq:initial-energy}.  Equal row-length multisets are the natural stronger invariant supplied by the skew Schur reduction and synchronize the two parents simultaneously for every admissible convex valuation.
\end{remark}

\section{Reconstruction from two boundary observations}\label{sec:reconstruction}

\subsection{The cofactor vector determines the diagram}

Vanishing of the defect is a statement about cofactors; being a translate is a statement about diagrams.  Proposition~\ref{prop:cofactor-initial} alone converts one into the other.

Write \(\rho_i=\lambda_i-\mu_i\) for the row lengths and
\[
o_i=\lambda_{i+1}-\mu_i
\qquad(0\le i\le n-2)
\]
for the adjacent overlaps, so that connectedness is \(o_i\ge1\).

\begin{lemma}\label{lem:distinguished}
For a connected \(X\) and \(0\le j\le n\), the initial monomial \eqref{eq:cofactor-initial} is, up to sign,
\begin{equation}\label{eq:Qj}
Q_j(X)
=
\prod_{1\le i<j}h_{o_{i-1}-1}\ \prod_{i\ge j}h_{\rho_i}.
\end{equation}
\end{lemma}

\begin{proof}
By \eqref{eq:Adef} and \eqref{eq:dexplicit}, \(d_{i+1}-A_i=\lambda_i-\mu_i=\rho_i\), while \(d_0-A_0=0\) and, for \(i\ge1\), \(d_i-A_i=\lambda_i-\mu_{i-1}-1=o_{i-1}-1\).  Substituting into \eqref{eq:cofactor-initial} gives \eqref{eq:Qj}, the factor \(h_0=1\) at \(i=0\) being omitted.
\end{proof}

The two extreme members are not new.  The subscripts of \(Q_0=\prod_i h_{\rho_i}\) are the row lengths and those of \(Q_n=\prod_i h_{o_{i-1}-1}\) are the adjacent overlaps less one, which are \cite[Prop.~6.2(ii) and (iii)]{Reiner2007Coincidences} respectively.  What \eqref{eq:Qj} adds is the interpolating family and the fact that consecutive members differ by one exchange, which is what the reconstruction below runs on.

Only two integers per cofactor are needed, rather than the cofactors themselves.  Here and in Proposition~\ref{prop:stable-cofactor} the valuation is the square one, \(\omega(k)=k^2\), which is what makes \(S_j\) a difference of squares and gives the closed form \eqref{eq:scalar-recovery}.  A general admissible weight still determines the pair, since \(z\mapsto\omega(z+\Delta)-\omega(z)\) is strictly increasing for each fixed \(\Delta\ge1\), but without a closed form.  Grade \(\Lam\) by \(\deg h_k=k\) and set
\begin{equation}\label{eq:scalar-data}
\delta_j=\deg C_j(X),
\qquad
\tau_j=\wt(C_j(X)),
\qquad
\Delta_j=\delta_j-\delta_{j+1},
\qquad
S_j=\tau_j-\tau_{j+1}.
\end{equation}

\begin{proposition}\label{prop:stable-cofactor}
Let \(X\) be connected with \(n\) nonempty rows.  Then \(\Delta_j>0\) and \(\Delta_j\) divides \(S_j\) for \(0\le j<n\), and
\begin{equation}\label{eq:scalar-recovery}
\rho_j
=
\frac12\left(\Delta_j+\frac{S_j}{\Delta_j}\right)
\quad(0\le j<n),
\qquad
o_{j-1}
=
1+\frac12\left(\frac{S_j}{\Delta_j}-\Delta_j\right)
\quad(1\le j<n).
\end{equation}
Consequently the ordered pairs \((\delta_j,\tau_j)_{j=0}^{n}\) determine \(X\) up to translation.  In particular two connected diagrams with equal cofactor vectors, \(C(A)=C(B)\), are translates of one another.
\end{proposition}

\begin{proof}
Each \(C_j(X)\) is homogeneous, being \(\pm s_{\lambda/\nu^{(j)}}\) by \eqref{eq:cofactor-schur}, and by Theorem~\ref{thm:monge-minor} its square-weight initial face is the single monomial \(\pm Q_j\) of \eqref{eq:Qj}.  So \(\delta_j\) and \(\tau_j\) are the first and second power sums of the index multiset of \(Q_j\), and the whole of \eqref{eq:scalar-data} is read off that multiset.

Read \eqref{eq:Qj} as a multiset of indices, with the index \(0\) omitted.  Passing from \(Q_j\) to \(Q_{j+1}\) removes one copy of \(\rho_j\) and inserts one copy of \(o_{j-1}-1\) when \(j\ge1\), and removes one copy of \(\rho_0\) and inserts nothing when \(j=0\).  Every row is nonempty, so \(\rho_j\ge1\) and the removal is visible.  For \(j\ge1\) the insertion cannot cancel the removal, because \(d_{j+1}-d_j=\mu_{j-1}-\mu_j+1=\rho_j-o_{j-1}+1\ge1\)
by Lemma~\ref{lem:grid-order}, so \(\rho_j\ge o_{j-1}>o_{j-1}-1\).  Hence the multiset difference \(Q_j\setminus Q_{j+1}\) is exactly \(\{\rho_j\}\) and \(Q_{j+1}\setminus Q_j\) is \(\{o_{j-1}-1\}\) or, when \(o_{j-1}=1\), empty.  Both \(\rho_j\) and \(o_{j-1}\) are thus determined.

Write \(z_0=0\) and \(z_j=o_{j-1}-1\) for \(j\ge1\), so that the exchange just described removes \(\rho_j\) and inserts \(z_j\).  Taking first and second power sums of that exchange gives
\[
\Delta_j=\rho_j-z_j,
\qquad
S_j=\rho_j^2-z_j^2=\Delta_j(\rho_j+z_j),
\]
and the previous paragraph gives \(\Delta_j>0\); hence \(\Delta_j\) divides \(S_j\), the two displayed linear equations in \(\rho_j\) and \(z_j\) are solvable, and their solution is \eqref{eq:scalar-recovery}.  The length of the sequence supplies \(n\).

Finally the profile determines the diagram up to translation: normalize \(\mu_{n-1}=0\), set \(\lambda_{n-1}=\rho_{n-1}\), and read upward by \(\mu_i=\lambda_{i+1}-o_i\) and \(\lambda_i=\mu_i+\rho_i\).  Equal cofactor vectors give equal \(Q_j\) and hence equal \eqref{eq:scalar-data}, so the last assertion is a special case.
\end{proof}

\subsection{Reconstruction from the first-row deletion}\label{subsec:child-separation}

\begin{lemma}\label{lem:top-row-cofactor}
For any \(X=\lambda/\mu\) with \(n\) nonempty rows, connected or not, and \(\bar X\) as above,
\[
C_1(X)=-s_{\bar X},
\]
with \(s_\varnothing=1\).
\end{lemma}

\begin{proof}
This is \eqref{eq:cofactor-schur} at \(j=1\), where the condition of \eqref{eq:cofactor-support} is empty, so connectedness is not consumed: the partition \(\nu^{(1)}\) of \eqref{eq:nu-j} is \((\lambda_0,\mu_1,\ldots,\mu_{n-1})\), so the top row of \(\lambda/\nu^{(1)}\) is empty and the diagram is \(\bar X\).  The sign is \((-1)^1\).
\end{proof}

\begin{lemma}\label{lem:two-cofactor-reconstruction}
For a connected skew diagram \(X\), the pair \((C_0(X),C_1(X))\) determines \(X\) up to translation.
\end{lemma}

\begin{proof}
Write \(\rho_i=d_{i+1}-A_i\), \(z_i=d_i-A_i\), with the monomials \(Q_j\) of \eqref{eq:Qj}, and fix any admissible weight \(\omega\).  The unique initial monomial of \(C_0\) is \(Q_0=\prod_ih_{\rho_i}\), so its number of factors, counted with multiplicity, is the number \(n\) of nonempty rows; every \(\rho_i\) is positive because every row is nonempty.  Also \(A_0=d_0=0\) and \(d_1=\rho_0=\deg C_0-\deg C_1\).  For \(n=1\) this already determines the diagram up to translation.

Suppose \(1\le k<n\) and that \(A_0,\ldots,A_{k-1}\) and \(d_0,\ldots,d_k\) are known.  Put
\[
T_j^{(k)}=\det M_{\{0,\ldots,k-1\},(2,\ldots,k,j)},
\]
the column tuple before \(j\) being empty when \(k=1\).  For \(2\le j\le k\) the column \(j\) is repeated, so \(T_j^{(k)}=0\) and only \(j\in\{0,1\}\) and \(j\ge k+1\) contribute.  This determinant is linear in its last column, so \((T_j^{(k)})_j\) is a linear combination of the first \(k\) rows of \(M\); since the signed cofactor vector lies in the right kernel of \(M\) (Section~\ref{sec:coordinates}), this gives
\[
R_k:=T_0^{(k)}C_0+T_1^{(k)}C_1=-\sum_{j=k+1}^{n}T_j^{(k)}C_j .
\]
The left-hand side is computable from the two observations and the recovered prefix.

Write \(\Pi_k=\prod_{i=0}^{k-2}h_{d_{i+2}-A_i}\).  For \(j\ge k+1\) the increasing matching in \(T_j^{(k)}\) sends row \(i<k-1\) to column \(i+2\) and row \(k-1\) to column \(j\); every matched column index is at least its row index, so the matching is feasible and Theorem~\ref{thm:monge-minor} gives \(\inin_\omega T_j^{(k)}=\Pi_kh_{d_j-A_{k-1}}\).  The successive differences of \(\wto(T_j^{(k)}C_j)\) are
\[
\omega(d_j-A_j)-\omega(d_{j+1}-A_j)+\omega(d_{j+1}-A_{k-1})-\omega(d_j-A_{k-1})>0
\]
by the strict Monge inequality, since \(A_j>A_{k-1}\).  Hence the summand at \(j=k+1\) alone carries the initial form, \(R_k\ne0\), and \(\inin_\omega R_k=(-1)^k\Pi_{k+1}Q_{k+1}\).

Let \(U_0=Q_1\) and let \(U_k\) be the unsigned initial monomial of \(R_k\).  With \(u_k=d_{k+1}-A_{k-1}\),
\[
\frac{U_k}{U_{k-1}}=\frac{h_{u_k}h_{z_k}}{h_{\rho_k}},
\qquad
u_k-\rho_k=A_k-A_{k-1}>0,
\qquad
\rho_k-z_k=d_{k+1}-d_k>0,
\]
and \(z_k\ge0\) by connectedness.  So \(u_k>\rho_k>z_k\ge0\): after cancellation the unique negative exponent is \(\rho_k\) and the smaller numerator index is \(z_k\), a single numerator factor meaning \(z_k=0\).  Therefore \(A_k=d_k-z_k\) and \(d_{k+1}=A_k+\rho_k\), and the induction closes after \(n-1\) steps.

Normalizing \(\mu_{n-1}=0\) gives \(L=d_n-n+1\), \(\lambda_i=L+i-A_i\) and \(\mu_i=L+i-d_{i+1}\), so both boundaries are determined.
\end{proof}

\begin{proof}[Proof of Theorem~\ref{thm:boundary-reconstruction}]
The first map supplies \(C_0=s_X\) and, by Lemma~\ref{lem:top-row-cofactor}, \(C_1=-s_{\bar X}\), so Lemma~\ref{lem:two-cofactor-reconstruction} applies.  For the second map apply that lemma to \(X\cdot(a)\), whose first-row deletion is \(X\).
\end{proof}

The determination in Theorem~\ref{thm:boundary-reconstruction} is constructive, and the construction is the induction of Lemma~\ref{lem:two-cofactor-reconstruction}: read \(n\) off the initial monomial of \(C_0\) and \(d_1\) from \(\deg C_0-\deg C_1\), then recover \(A_k\) and \(d_{k+1}\) in turn from the initial monomials of the prefix residuals \(R_k\), and normalize \(\mu_{n-1}=0\).  Enumerating assignments of the row lengths would not do, because the inner boundary is not supplied: the connected diagrams \((3+t,3)/(t,0)\) for \(t=0,1,2\) share the row-length multiset \(\{3,3\}\), and hence the initial monomial \(h_3^2\) of \(C_0\), while their inner boundaries differ.

\subsection{Reconstruction from two positive extensions}\label{subsec:two-positive-extensions}

\begin{proof}[Proof of Corollary~\ref{cor:two-positive-extensions}]
Let \(X\) be connected, with no boundary prescribed.  Grade \(\Lam\) by \(\deg h_k=k\).  Since \(\deg s_{X\cdot(k)}=|X|+k\) and the two lengths are distinct, the two observations have distinct degrees; order them as \(F,H\) with \(\deg F<\deg H\) and put \(c=\deg H-\deg F>0\).  Writing \(a<b\) for the unknown lengths, \(F=s_{X\cdot(a)}\), \(H=s_{X\cdot(b)}\) and \(c=b-a\).  Put \(Z=(X\cdot(a))\cdot(c)\); every attachment has positive length and meets the row below it in exactly one column, so \(Z\) is connected.  By \eqref{eq:iterated-cap} the two observations determine
\[
s_Z=h_cs_{X\cdot(a)}-s_{X\cdot(b)},
\qquad
s_{\bar Z}=s_{X\cdot(a)},
\]
both computable from the observations alone.  Theorem~\ref{thm:boundary-reconstruction} applied to \(Z\) determines it up to translation.  Deleting its first two rows recovers \(X\), and its second row has length \(a\), so \(b=a+c\).
\end{proof}

\begin{remark}\label{rem:one-extension-insufficient}
Two extensions are needed, and not merely at one exceptional length: for every
\(a\ge1\) a single observation \(s_{X\cdot(a)}\) fails to determine \(X\).  With
\[
\mu^{[a]}=(4a-2,\,2a-1,\,0),
\qquad
U_a=(6a-2,\,5a-1,\,2a)/\mu^{[a]},
\qquad
V_a=(6a-2,\,4a-1,\,3a)/\mu^{[a]},
\]
both parents are connected, and
\[
s_{U_a\cdot(a)}=s_{V_a\cdot(a)}
\qquad\text{while}\qquad
s_{U_a}-s_{V_a}=h_{3a}h_{4a}-h_{2a}h_{5a}\ne0 .
\]
The verification is exact rather than a search.  At \(a=1\) the two extensions
are the equal pair of Example~\ref{ex:actual-equal-chain} below, and the general
member is its image under the substitution \(h_k\mapsto h_{ak}\): every
Jacobi--Trudi subscript of \(U_a,V_a\) and of their extensions is \(a\) times
the corresponding subscript at \(a=1\), and that substitution is an injective
ring homomorphism of \(\Lam\), so it carries the equality to an equality and the
nonzero difference to a nonzero difference.  What this shows is that one
observation is insufficient in general at each fixed length; it does not say
that every prescribed inner boundary requires two.
\end{remark}



\section{First-split localization and cap-fixed rigidity}\label{sec:localization}

Assume in this section that the two parents have a common cap grid and equal row-length multisets, except in Theorem~\ref{thm:cap-fixed-rigidity}, whose floor hypothesis supplies the multisets rather than assuming them.  The first changed row \(v\) and the barrier \(q=v+1\) are those of \eqref{eq:first-change} and \eqref{eq:q-def}.

This section and the two that follow concern a \emph{pair} of diagrams, where Section~\ref{sec:reconstruction} concerned one at a time.

\subsection{First-split localization}\label{subsec:first-split}

\begin{proposition}\label{prop:synchronization}
For every admissible convex weight \(\omega\),
\begin{equation}\label{eq:synchronized}
\sigma_j^{A,\omega}=\sigma_j^{B,\omega}
\qquad(0\le j\le v).
\end{equation}
At the next column,
\begin{equation}\label{eq:first-split}
\sigma_q^{A,\omega}\ne\sigma_q^{B,\omega}.
\end{equation}
The direction of the inequality in \eqref{eq:first-split} is determined by the sign of \(A_v^A-A_v^B\) and is independent of \(\omega\).
\end{proposition}

\begin{proof}
Equality at \(j=0\) follows from Proposition~\ref{prop:initial-energy}.  The action increments coincide for \(j<v\) because the corresponding row parameters coincide, giving \eqref{eq:synchronized}.  At \(j=v\), Proposition~\ref{prop:action-recurrence} and strict monotonicity in \(A_v\) give a strict split.
\end{proof}

\begin{corollary}\label{cor:first-split-anchor}
For every admissible convex weight,
\[
G_q\ne0
\]
and
\[
\alpha_q^\omega
=
\min\{\sigma_q^{A,\omega},\sigma_q^{B,\omega}\}.
\]
\end{corollary}

\begin{proof}
The two parent cofactor valuations at \(q\) are different, so by Lemma~\ref{lem:valuation-laws} their difference cannot cancel at the lower valuation.
\end{proof}

\subsection{Common-row relations and tropical balancing}\label{subsec:balancing}

The relations \eqref{eq:common-row-relation} are exact, which restricts where their lowest-valuation summand can sit.

\begin{lemma}\label{lem:tropical-balancing}
For every \(i<v\), the minimum of \(F_i^\omega\) is finite and is attained at least twice.
\end{lemma}

\begin{proof}
The minimum is finite because Corollary~\ref{cor:first-split-anchor} gives \(G_q\ne0\) and Lemma~\ref{lem:diagonal-feasible} gives \(d_q\ge A_i\), so the summand at \(q\) has finite tilted cost.  If a single summand had uniquely smallest valuation in \eqref{eq:common-row-relation}, then by Lemma~\ref{lem:valuation-laws} the initial form of the entire sum would equal the nonzero initial form of that summand, contradicting the exact zero relation.
\end{proof}

Geometrically, for square weight the points \(p_j=(d_j,\alpha_j)\) are exposed by supporting lines of slope \(2A_i\).  Lemma~\ref{lem:tropical-balancing} says each common row must expose a nontrivial lower-hull face rather than a single point.  Such a face spanning two consecutive columns is an \emph{adjacent edge}, and a minimizer set \(\mathcal M_i=\{i+1,i+2\}\) is the adjacent edge at row \(i\).

\subsection{No-backflow: the first changed row as a causal barrier}\label{subsec:no-backflow}

For a common-prefix row \(i<v\), define the tilted parent quantities
\begin{equation}\label{eq:F-parent}
F_i^{X,\omega}(j)
=
\sigma_j^{X,\omega}
+
\omega(d_j-A_i)-\omega(d_j).
\end{equation}

\begin{lemma}\label{lem:tail-increase}
If \(k\ge q=v+1\), then
\[
F_i^{X,\omega}(k+1)>F_i^{X,\omega}(k)
\qquad(i<v).
\]
\end{lemma}

\begin{proof}
Combining Proposition~\ref{prop:action-recurrence} with \eqref{eq:F-parent}, the increment is
\[
\begin{aligned}
&\omega(d_k-A_k^X)+\omega(d_{k+1}-A_i)\\
&\qquad-
\omega(d_{k+1}-A_k^X)-\omega(d_k-A_i).
\end{aligned}
\]
Here \(A_i<A_k^X\) and \(d_k<d_{k+1}\), so the Strict-Monge inequality makes the expression positive.
\end{proof}

\begin{theorem}[No-backflow]\label{thm:no-backflow}
For every common-prefix row \(i<v\) and every \(\ell>q\),
\begin{equation}\label{eq:no-backflow}
F_i^\omega(\ell)>F_i^\omega(q).
\end{equation}
For the square weight the separation is quantitative:
\begin{equation}\label{eq:no-backflow-margin}
F_i(\ell)-F_i(q)
\ge
(\ell-q)(\ell+q-2i-1)
\qquad(q<\ell\le n,\ i<v).
\end{equation}
\end{theorem}

\begin{proof}
For every defect coordinate,
\[
\alpha_\ell^\omega
\ge
\min_X\sigma_\ell^{X,\omega}.
\]
Thus
\[
F_i^\omega(\ell)
\ge
\min_X F_i^{X,\omega}(\ell).
\]
Lemma~\ref{lem:tail-increase} makes each parent tilted profile strictly increase beyond \(q\), while Corollary~\ref{cor:first-split-anchor} gives equality with the smaller parent profile at \(q\).  Therefore \eqref{eq:no-backflow} follows.

For \eqref{eq:no-backflow-margin}, the increment computed in Lemma~\ref{lem:tail-increase} is \(2(A_k^X-A_i)(d_{k+1}-d_k)\) at the square weight, so summing it from \(q\) to \(\ell-1\) and taking the smaller parent at each step gives
\[
F_i(\ell)-F_i(q)
\ge
2\sum_{k=q}^{\ell-1}\bigl(\min\{A_k^A,A_k^B\}-A_i\bigr)(d_{k+1}-d_k).
\]
Lemma~\ref{lem:grid-order} makes both parameter sequences strictly increasing integers, so \(\min\{A_k^A,A_k^B\}-A_i\ge k-i\) and \(d_{k+1}-d_k\ge1\), whence the sum is at least \(2\sum_{k=q}^{\ell-1}(k-i)=(\ell-q)(\ell+q-2i-1)\).  The bound measures separation from the anchor and is not claimed sharp, nor does it assert that the later defect coordinates increase monotonically.
\end{proof}

\begin{remark}
Theorem~\ref{thm:no-backflow} explains why arbitrarily many later changed rows do not obstruct the chain.  Once the first changed row occurs, nothing later in the diagram can propagate backward and determine the leading cancellation equations of the common prefix.  Later rows may alter coefficients, supports, or arbitrarily deep residual strata, but they lie outside the causal cone relevant to Theorem~\ref{thm:floor}.
\end{remark}

\subsection{Marching minima}\label{subsec:marching}

Let \(R_i^\omega=\max\mathcal M_i^\omega\).
By Lemma~\ref{lem:tropical-balancing}, \(|\mathcal M_i^\omega|\ge2\),
and by Theorem~\ref{thm:no-backflow},
\begin{equation}\label{eq:R-upper}
R_i^\omega\le q.
\end{equation}

\begin{lemma}\label{lem:marching}
For \(0\le i<v-1\),
\[
R_{i+1}^\omega>R_i^\omega,
\]
and the minimizer sets interlace:
\[
\min\mathcal M_{i+1}^\omega\ge R_i^\omega=\max\mathcal M_i^\omega.
\]
\end{lemma}

\begin{proof}
Put \(r=R_i^\omega\).  If \(d_r<A_{i+1}\), then every column \(j\le r\) disappears from the next row, and the result is immediate.

Otherwise \(r\) remains available.  For any available \(j<r\), \(F_i^\omega(j)\ge F_i^\omega(r)\).
Strict-Monge single crossing implies
\[
F_{i+1}^\omega(j)-F_{i+1}^\omega(r)
>
F_i^\omega(j)-F_i^\omega(r)
\ge0.
\]
Hence nothing left of \(r\) minimizes in the next row, and a minimizer exists because Lemma~\ref{lem:tropical-balancing} gives a finite minimum in row \(i+1\) as well.  If \(r\) itself remains minimal, tropical balancing forces a second minimum to its right; if not, every new minimizer is to its right.  Thus \(R_{i+1}^\omega>r\).
\end{proof}

\subsection{Adjacent-edge rigidity at a fixed cap}

\begin{theorem}[Cap-fixed adjacent-edge rigidity]\label{thm:cap-fixed-rigidity}
Assume connected parents with a common cap grid and first changed row \(v\ge1\), with \(G\) their cofactor defect \eqref{eq:defect}.  Equal row-length multisets are \emph{not} assumed; the floor hypothesis supplies them.  If
\begin{equation}\label{eq:floor-antecedent-general}
\alpha_0^\omega>m_+^\omega,
\end{equation}
then
\begin{equation}\label{eq:adjacent-chain-capfixed}
\mathcal M_i^\omega=\{i+1,i+2\}
\qquad(0\le i<v).
\end{equation}
Consequently
\begin{equation}\label{eq:floor-set-general}
T_{0,\omega}^+=\{d_1,d_2\},
\end{equation}
\begin{equation}\label{eq:support-prefix}
G_j\ne0\qquad(1\le j\le v+1),
\end{equation}
and
\begin{equation}\label{eq:prefix-profile}
\alpha_1^\omega=\alpha_2^\omega<\alpha_3^\omega<\cdots<\alpha_{v+1}^\omega.
\end{equation}
Every later positive coordinate has normalized action strictly larger than \(\alpha_{v+1}^\omega\).
\end{theorem}

\begin{proof}
The floor hypothesis supplies the initial energies.  Write \(\mathcal E_X=\sigma_0^{X,\omega}\), and put \(\mathcal E=\min(\mathcal E_A,\mathcal E_B)\).  Since \(A_0=0\) by \eqref{eq:Adef}, the action recurrence gives \(\sigma_1^{X,\omega}=\sigma_0^{X,\omega}\), and \(A_j>0\) for \(j\ge1\) makes the later increments positive, so \(\sigma_j^{X,\omega}\ge\mathcal E_X\) for every \(j\).  Hence \(\alpha_j^\omega\ge\mathcal E\) for every \(j\), and in particular \(m_+^\omega\ge\mathcal E\).

By \eqref{eq:floor-antecedent-general} we get \(\wto(G_0)=\alpha_0^\omega>\mathcal E\).  Were \(\mathcal E_A\ne\mathcal E_B\), the initial form of the parent of smaller energy would survive in \(G_0\) and give \(\wto(G_0)=\mathcal E\); so \(\mathcal E_A=\mathcal E_B\), and their initial forms cancel.  Each is a single monomial with coefficient \(1\) by Proposition~\ref{prop:cofactor-initial}, so the two coincide, and by \eqref{eq:initial-energy} they are the row-length products.  The two parents therefore have equal initial \(\omega\)-energy and equal row-length multisets, which is what Proposition~\ref{prop:synchronization} consumes.

At row \(0\), column \(0\) is not minimizing by \eqref{eq:floor-antecedent-general}.  Thus \(R_0^\omega\ge2\).  After \(v-1\) strict marching steps, \(R_{v-1}^\omega\ge R_0^\omega+v-1\ge v+1=q\).
The upper bound \eqref{eq:R-upper} forces equality throughout.  Hence \(R_0^\omega=2\) and \(R_i^\omega=i+2\).  Since column \(0\) does not minimize and each row has at least two minimizers, row \(0\) must have exactly \(\{1,2\}\).  The interlacing of Lemma~\ref{lem:marching} then confines \(\mathcal M_i^\omega\) to \(\{i+1,i+2\}\), where two minimizers force equality.  The rightmost values alone would not do this: they leave an extra minimizer far to the left unexcluded.

The equality \(\mathcal M_0^\omega=\{1,2\}\) gives \eqref{eq:floor-set-general}.  Every adjacent tie in \eqref{eq:adjacent-chain-capfixed} forces both coordinates to be supported, giving \eqref{eq:support-prefix}.  The exact valuation differences follow from equality of the two adjacent tilted costs.  For \(i\ge1\), strict convexity makes these differences positive.  Finally, the parent action profiles increase strictly beyond the first-split anchor, and defect cancellation can only raise valuation, proving the tail statement.
\end{proof}

The exact adjacent recurrence in Theorem~\ref{thm:cap-fixed-rigidity} is
\[
\begin{aligned}
\alpha_{i+2}^\omega-\alpha_{i+1}^\omega
&=
\omega(d_{i+2})-\omega(d_{i+2}-A_i)\\
&\quad-
\omega(d_{i+1})+\omega(d_{i+1}-A_i),
\end{aligned}
\]
for \(0\le i<v\).  At \(i=0\) the right-hand side vanishes, yielding \(\alpha_1^\omega=\alpha_2^\omega\).

\begin{remark}
The theorem implies
\[
\{d_1,\ldots,d_{v+1}\}\subseteq\supp G.
\]
Thus if \(m\) denotes the number of positive supported shifts, then the floor antecedent forces
\[
m\ge v+1.
\]
This is a consequence of the Monge geometry rather than an input to it.
\end{remark}

\subsection{Application to actual skew Schur equality}\label{subsec:actual-equality}

For an actual equal pair the floor antecedent is automatic, because the zeroth defect coordinate is the difference of the two skew Schur functions.  The cap grid is automatic as well.

The invariant is implicit in the proof of \cite[Cor.~6.3, eq.~(6.1)]{Reiner2007Coincidences}, where the coefficient of the unique top-right generator is a complementary minor shown to be nonzero; the irreducibility argument that proof serves was completed in \cite{Barekat2009CoincidencesCorrigendum}.  We record the invariant with the translation term \(\mu_{n-1}\) explicit and give the non-vanishing argument in the present coordinates, since equality uses it to align the cap grids.

\begin{lemma}\label{lem:equality-normalizes-cap}
Let \(X=\lambda/\mu\) be connected with \(n\) nonempty rows.  Then the largest generator index occurring in \(s_X\) is
\begin{equation}\label{eq:largest-index}
H_X=\lambda_0-\mu_{n-1}+n-1.
\end{equation}
If \(Y=\kappa/\eta\) is connected and \(s_X=s_Y\), then:
\begin{enumerate}[label=(\roman*),leftmargin=2em]
\item the row-length multisets of \(X\) and \(Y\) agree, and \(Y\) has \(n\) nonempty rows;
\item if moreover \(\eta=\mu\), then \(\lambda_0=\kappa_0\).
\end{enumerate}
\end{lemma}

\begin{proof}
The increasing matching of the ordinary Jacobi--Trudi matrix has monomial \(\prod_{i=0}^{n-1}h_{\lambda_i-\mu_i}\), which by Theorem~\ref{thm:monge-minor} is its unique initial monomial and has coefficient one.  Equality of the polynomials therefore gives equality of the row-length multisets; every row length of a connected diagram is positive, so the two multisets have the same number of members and \(Y\) has \(n\) nonempty rows.

The largest subscript in that matrix is \(H_X\) of \eqref{eq:largest-index}, and it occurs only in the top-right entry, since \(\lambda_i-i\) is strictly decreasing and \(j-\mu_j\) strictly increasing.  Its coefficient in \(s_X\) is therefore, up to sign, the minor deleting the first row and the last column.  The increasing matching of that minor has subscripts \(\lambda_i-\mu_{i-1}-1\) for \(1\le i<n\), which are nonnegative by connectedness, so Theorem~\ref{thm:monge-minor} makes the minor nonzero; for \(n=1\) it is the empty determinant.  This proves \eqref{eq:largest-index}: \(H_X\) is the largest generator index that actually occurs, not merely the largest available in the presentation.  Equality gives \(H_X=H_Y\).  If moreover \(\eta=\mu\), the common inner boundary and row count give \(\lambda_0=\kappa_0\).
\end{proof}

\begin{theorem}[Rigidity of actual equal pairs]\label{thm:equal-pair-rigidity}
Let \(X=\lambda/\mu\) and \(Y=\kappa/\mu\) be distinct connected skew diagrams with \(n\) nonempty rows and \(s_X=s_Y\).  Then \(\lambda_0=\kappa_0\), so their cap grids coincide.  Let \(v\) be the first row at which the outer boundaries differ and put \(G=C(X)-C(Y)\).  For every admissible convex weight \(\omega\),
\begin{equation}\label{eq:equal-pair-chain}
G_0=0,
\qquad
\mathcal M_i^\omega=\{i+1,i+2\}
\quad(0\le i<v),
\qquad
T_{0,\omega}^+=\{d_1,d_2\}.
\end{equation}
In particular \(v\ge1\), and \(G_j\ne0\) for \(1\le j\le v+1\).
\end{theorem}

\begin{proof}
Lemma~\ref{lem:equality-normalizes-cap} puts the parents on a common cap grid and, through Proposition~\ref{prop:initial-energy}, supplies equal initial \(\omega\)-energy for every admissible weight.  Their top outer rows agree, so the first changed row satisfies \(v\ge1\).  Corollary~\ref{cor:first-split-anchor} gives \(G_{v+1}\ne0\); hence \(m_+^\omega\) is finite, and in particular \(G\ne0\) without appeal to Proposition~\ref{prop:stable-cofactor}.  But \(G_0=C_0(X)-C_0(Y)=s_X-s_Y=0\),
so \(\alpha_0^\omega=+\infty>m_+^\omega\).  Theorem~\ref{thm:cap-fixed-rigidity} applies directly and gives \eqref{eq:equal-pair-chain}.
\end{proof}

\begin{remark}\label{rem:no-virtual-row}
The chain here starts at \(\mathcal M_0^\omega\), not at the virtual row \(\mathcal M_{-1}^\omega=\{0,1\}\) of Theorem~\ref{thm:ratio-rigidity}.  The virtual edge belongs to the nonzero-defect ratio problem, where \(G_0\ne0\); on an actual equal pair the zeroth coordinate vanishes and the antecedent comes for free instead.
\end{remark}

\begin{proposition}\label{prop:monomial-prefix}
In the situation of Theorem~\ref{thm:equal-pair-rigidity} put \(q=v+1\).  There are \(h\)-monomials \(M_1,\ldots,M_q\) and signs \(\varepsilon_1,\ldots,\varepsilon_q\), all independent of \(\omega\), with
\[
\inin_\omega G_j=\varepsilon_jM_j
\quad(1\le j\le q),
\qquad
\varepsilon_j=-\varepsilon_{j+1}
\quad(1\le j<q),
\]
and
\begin{equation}\label{eq:monomial-prefix-exchange}
h_{d_j-A_{j-1}}M_j
=
h_{d_{j+1}-A_{j-1}}M_{j+1}
\qquad(1\le j<q).
\end{equation}
\end{proposition}

\begin{proof}
At \(q\) the two parent cofactor valuations differ, with the same direction for every admissible weight by Proposition~\ref{prop:synchronization}, so the lower parent alone contributes the initial form of \(G_q\); its initial monomial and unit coefficient are those of Proposition~\ref{prop:cofactor-initial} and do not depend on \(\omega\).

For \(1\le j<q\) the exact common-row relation \eqref{eq:common-row-relation} at row \(j-1\) has exactly the two minimizing columns \(j\) and \(j+1\) by \eqref{eq:equal-pair-chain}, so its initial form is
\[
h_{d_j-A_{j-1}}\inin_\omega G_j
+
h_{d_{j+1}-A_{j-1}}\inin_\omega G_{j+1}
=0.
\]
Multiplication by a monomial in \(\Lam\) is injective and preserves the number of terms, so descending from \(q\) keeps each initial form a single monomial, with coefficient \(-\varepsilon_{j+1}\), determined by \eqref{eq:monomial-prefix-exchange} from \(M_{j+1}\) and the common grid alone.
\end{proof}

\subsection{An actual equal pair}\label{subsec:equal-example}

The pair below is a smallest instance of Theorem~\ref{thm:equal-pair-rigidity}.

\begin{example}\label{ex:actual-equal-chain}
Let
\[
X=(4,4,4,2)/(3,2,1,0),
\qquad
Y=(4,4,3,3)/(3,2,1,0).
\]
Their common cap grid is \((0,1,3,5,7)\) and their row-parameter sequences are \((0,1,2,5)\) and \((0,1,3,4)\).  Expanding the maximal minors gives
\[
G_0=G_4=0,
\qquad
G_1=h_2h_5-h_3h_4,
\qquad
G_2=h_1h_4-h_5,
\qquad
G_3=h_3-h_1h_2,
\]
so \(s_X=s_Y\) and the first changed row is \(v=2\).  The two common-row relations \eqref{eq:common-row-relation} are
\[
h_1G_1+h_3G_2+h_5G_3=0,
\qquad
G_1+h_2G_2+h_4G_3=0.
\]
For square weight the normalized actions are
\[
(\alpha_0,\alpha_1,\alpha_2,\alpha_3,\alpha_4)
=(+\infty,26,26,30,+\infty),
\]
whence \(\mathcal M_0=\{1,2\}\) and \(\mathcal M_1=\{2,3\}\).  The first-row deletions are separated by
\[
s_{\bar X}-s_{\bar Y}=h_3h_4-h_2h_5\ne0,
\]
which is Theorem~\ref{thm:boundary-reconstruction} on this pair.
Figure~\ref{fig:cancellation-geometry} shows the square-weight geometry of
the same pair.
\end{example}

\begin{figure}[htbp]
\centering
\includegraphics[width=\linewidth]{figures/monge-cancellation-geometry.pdf}
\caption{Square-weight geometry of Example~\ref{ex:actual-equal-chain}.  The
finite points are \((d_j,\alpha_j)\) for \(j=1,2,3\); \(G_0=G_4=0\), so those
coordinates have infinite action and contribute no point.  A common row \(i\)
minimizes \(F_i\) of \eqref{eq:F-general} along lines of slope \(2A_i\), so the
two common rows expose the adjacent edges \(\mathcal M_0=\{1,2\}\) and
\(\mathcal M_1=\{2,3\}\) --- the chain of
Theorem~\ref{thm:equal-pair-rigidity} on this pair.  A positive cap \(a\)
weights the summand at column \(j\) by \(\alpha_j-\omega(d_j)+\omega(a+d_j)\),
whose level sets have slope \(-2a\); at \(a=1\) the support touches column one
alone, which is Corollary~\ref{cor:boundary-separation}.  This is one weight,
not every admissible convex weight, and the negative-slope support is not a
virtual-row relation: \(G_0\) vanishes here
(Remark~\ref{rem:no-virtual-row}).}
\label{fig:cancellation-geometry}
\end{figure}

\section{Ratio identities and full adjacent-edge rigidity}\label{sec:full-rigidity}

In this section the parents are connected, have a common inner boundary,
and their nonzero cofactor defect satisfies \eqref{eq:ratio-identity}.
Neither equal row-length multisets nor a common cap grid is assumed: the
identity supplies the row data first, then the grid, and then one further
row parameter to the left of the genuine top row.

\begin{lemma}\label{lem:ratio-row-lengths}
For connected parents, a ratio identity forces equal row-length multisets,
and hence the same number of nonempty rows.  No common inner boundary is
needed for this conclusion.
\end{lemma}

\begin{proof}
At \(r=0\), \eqref{eq:P0-vanishes} makes the identity read
\(s_{A\cdot(a)}=s_{B\cdot(a)}\); at \(r>0\), Lemma~\ref{lem:extension-cofactors}
makes it read \(s_{(A\cdot(r))\cdot(a)}=s_{(B\cdot(r))\cdot(a)}\).  Either way the two
sides are equal skew Schur functions of connected diagrams, since \(X\cdot(k)\)
meets the top row of \(X\) in exactly one column and an extension of a connected
diagram is therefore connected.
Lemma~\ref{lem:equality-normalizes-cap}(i) equates their row-length multisets.
Extending by a row of length \(k\) adjoins \(k\) to the multiset and changes
nothing else, so removing the adjoined \(a\), or \(a\) and \(r\), leaves
\(\rho(A)=\rho(B)\).  Every parent row is nonempty, so the cardinalities give
the common row count.
\end{proof}

\subsection{Boundary normalization}\label{subsec:ratio}

Let \(D\) be the largest cap shift occurring in either parent.  Since every cofactor entry has the form \(h_{d_j-A_i}\) with \(A_i\ge0\), every cofactor coefficient belongs to
\begin{equation}\label{eq:bounded-generator}
\Z[h_1,\ldots,h_D].
\end{equation}

Let \(e=\max\supp G\).

\begin{theorem}\label{thm:terminal-gap}
Assume connected parents with a common inner boundary and nonzero defect, and assume \eqref{eq:ratio-identity}.  Then
\[
e+r+a\le D.
\]
In particular \(e<D\).
\end{theorem}

\begin{proof}
If \(H=e+r+a>D\), regard \(P_{r+a}(G)-h_aP_r(G)=0\)
as a polynomial identity in the free generator \(h_H\).  No \(G_d\) contains \(h_H\) by \eqref{eq:bounded-generator}.  The external factor \(h_{r+a+d}\) equals \(h_H\) only at \(d=e\).  On the right, \(h_{r+d}=h_H\) would require \(d=e+a>e\).  The remaining occurrence is the multiplier \(h_a\) itself, which is \(h_H\) exactly when \(e=r=0\); that case is excluded because \eqref{eq:P0-vanishes} gives \(P_0(G)=0\).  At \(e=0\) the support is \(\{0\}\), so \eqref{eq:Pk} makes that relation read \(G_0h_0=G_0=0\), hence \(G=0\).  We may therefore assume \(e\ge1\), and the coefficient of \(h_H\) is exactly \(G_e\ne0\), contradiction.
\end{proof}

\begin{remark}\label{rem:realizability-terminal}
The vector \(G_0=1\), \(G_d=0\) for \(d>0\), with \(r=0\), satisfies \eqref{eq:ratio-identity} for every \(a\), and for it \(e+r+a=a\) exceeds \(D\) as soon as \(a\) does.  Theorem~\ref{thm:terminal-gap} is therefore false for an arbitrary vector and true for a defect, and what separates the two is \eqref{eq:P0-vanishes}: a cofactor vector of an actual Jacobi--Trudi board satisfies its own row relations.  The same distinction recurs at Theorem~\ref{thm:floor}; see Section~\ref{sec:scope}.
\end{remark}

\begin{proposition}\label{prop:top-terminal}
Suppose the two connected parents have common \(\mu\), the same number of nonempty rows, and different top outer rows.  Then the larger-top-row parent has a unique largest cap shift, and the defect is nonzero at that shift.
\end{proposition}

\begin{proof}
Orient the pair so that \(\lambda_0^B=\lambda_0^A+\theta\) with \(\theta>0\).  Then \(b_0\) increases by \(\theta\), while every non-cap \(b_j\) remains common.  Hence every nonzero cap shift of \(B\) increases by \(\theta\), and its terminal shift is strictly larger than every shift of \(A\).  The corresponding terminal maximal cofactor is nonzero by Theorem~\ref{thm:monge-minor}; hence the defect at that unique largest shift is nonzero.
\end{proof}

\begin{corollary}\label{cor:cap-fixing}
Under \eqref{eq:ratio-identity},
\[
\lambda_0^A=\lambda_0^B.
\]
With common \(\mu\), the two parents therefore have a common cap grid.
\end{corollary}

\begin{proof}
Lemma~\ref{lem:ratio-row-lengths} gives the two parents the same number of nonempty rows, which is the remaining hypothesis of Proposition~\ref{prop:top-terminal}.  Were their top outer rows different, that proposition would give \(e=D\), contradicting Theorem~\ref{thm:terminal-gap}.
\end{proof}

\begin{remark}
The ratio identity does not create the interior Strict-Monge rigidity.  It forces both diagrams into the same ordered cap coordinate system.  Once the cap is fixed, Theorem~\ref{thm:cap-fixed-rigidity} no longer uses the ratio identity.
\end{remark}

\subsection{The ratio identity as a virtual row}\label{subsec:virtual}

The ratio identity also contributes one additional row parameter to the left of the genuine top row.

Define
\begin{equation}\label{eq:gamma}
\gamma=
\begin{cases}
 a,&r=0,\\
 r,&r>0.
\end{cases}
\end{equation}

\subsubsection{The case \(r=0\)}

Since the common top row gives \(P_0(G)=0\),
the ratio identity becomes \(P_a(G)=0\).
Its \(j\)-th summand has valuation
\[
\wto(G_j)+\omega(a+d_j)
=
\alpha_j^\omega+\omega(a+d_j)-\omega(d_j).
\]
This is exactly the tilted cost corresponding to the virtual parameter \(A_{-1}=-a\).

\subsubsection{The case \(r>0\)}

Write \(Q_j=h_{r+a+d_j}-h_a h_{r+d_j}\).
Because \(\omega\) has positive increasing increments, it is strictly superadditive on positive integers:
\[
\omega(x+y)>\omega(x)+\omega(y)
\qquad(x,y>0).
\]
Therefore \(\omega(r+a+d_j) > \omega(a)+\omega(r+d_j)\),
so
\[
\inin_{\omega}(Q_j)=-h_a h_{r+d_j}
\]
and \(\wto(Q_j)=\omega(a)+\omega(r+d_j)\).
Up to the common additive constant \(\omega(a)\), the ratio summand therefore has tilted cost \(\alpha_j^\omega+\omega(r+d_j)-\omega(d_j)\),
which is the virtual row \(A_{-1}=-r\).

Thus in both cases the ordered row parameters begin \(-\gamma=A_{-1}<A_0=0<A_1<\cdots<A_{v-1}\).

The virtual index carries the hypotheses its relation needs, so the results of Section~\ref{sec:localization} apply at \(i=-1\) as they do at a genuine common row.  Lemma~\ref{lem:tropical-balancing}, Theorem~\ref{thm:no-backflow} and Lemma~\ref{lem:marching} are stated for common-prefix rows, and each consumes only exactness of the relation at that index together with \(A_i<A_k^X\) for every genuine row parameter, which is the display above; \eqref{eq:F-general} supplies the tilted cost at \(A_{-1}=-\gamma\), and Lemma~\ref{lem:strict-monge} is stated on differences, so \(A_{-1}<0\) costs nothing.  This is how the chain below reaches one step further left than the cap-fixed one.

\subsubsection{The complete minimizing chain}

\begin{theorem}\label{thm:ratio-rigidity}
Assume connected parents with common \(\mu\), nonzero defect, and the ratio identity \eqref{eq:ratio-identity}, with \(G\) their cofactor defect \eqref{eq:defect}.  Let \(v\) be the first changed row after Corollary~\ref{cor:cap-fixing}.  Then for every admissible convex weight \(\omega\) the tropical minimizer sets are
\begin{equation}\label{eq:virtual-edge}
\mathcal M_{-1}^\omega=\{0,1\},
\end{equation}
and
\begin{equation}\label{eq:full-edge-chain}
\mathcal M_i^\omega=\{i+1,i+2\}
\qquad(0\le i<v).
\end{equation}
Hence
\begin{equation}\label{eq:universal-floor-set}
T_{0,\omega}^+=\{d_1,d_2\}
\end{equation}
and the cap-fixed floor antecedent is automatic:
\begin{equation}\label{eq:automatic-floor-antecedent}
\alpha_0^\omega>m_+^\omega.
\end{equation}
\end{theorem}

\begin{proof}
The virtual relation satisfies the same tropical balancing principle as a genuine common row, and no-backflow holds because \(A_{-1}<A_k^X\) for every genuine row parameter.  Thus its rightmost minimizer \(R_{-1}^\omega\) is at least \(1\).  There are now \(v\) strict marching steps before row \(v-1\), while no-backflow still gives \(R_{v-1}^\omega\le q=v+1\).  Hence \(R_{v-1}^\omega\ge R_{-1}^\omega+v\ge v+1=q\),
so equality holds at every step.  Therefore \(R_{-1}^\omega=1\) and \(R_i^\omega=i+2\).  Interlacing and tropical balancing together force the exact adjacent pairs in \eqref{eq:virtual-edge}--\eqref{eq:full-edge-chain}, as in Theorem~\ref{thm:cap-fixed-rigidity}.  In particular row \(0\) has exactly the positive minima \(\{1,2\}\), yielding \eqref{eq:universal-floor-set} and \eqref{eq:automatic-floor-antecedent}.
\end{proof}

\subsubsection{Exact gap formulas}

The virtual tie gives
\begin{equation}\label{eq:virtual-gap-general}
\alpha_0^\omega-\alpha_1^\omega
=
\omega(d_1+\gamma)-\omega(d_1)-\omega(\gamma)+\omega(0).
\end{equation}
Strict convexity makes this quantity positive.

Each genuine adjacent tie gives the recurrence displayed after Theorem~\ref{thm:cap-fixed-rigidity}, so the profile \eqref{eq:prefix-profile} holds here as well, now with the virtual gap \(\alpha_0^\omega>\alpha_1^\omega\) in front of it and every later positive coordinate higher still.

For square weight, \eqref{eq:virtual-gap-general} becomes
\begin{equation}\label{eq:square-gap}
\alpha_0-m_+
=2\gamma d_1
=\begin{cases}
2a d_1,&r=0,\\
2r d_1,&r>0.
\end{cases}
\end{equation}

\subsection{Cancellative Leftmost as a square-weight specialization}\label{subsec:floor}

We now specialize to the square valuation \(\omega(k)=k^2\).
Write \(\alpha_j=\wt(G_j)+d_j^2\),
and \(T_0^+=\{d_j:j>0,\ \alpha_j=m_+\}\).
Let \(t\) be the leftmost positive supported shift.

The automatic top-row relation is
\begin{equation}\label{eq:P0}
P_0(G)=G_0+\sum_{j>0}h_{d_j}G_j=0.
\end{equation}
Define
\begin{equation}\label{eq:J0}
J_0
=
\sum_{d_j\in T_0^+}h_{d_j}\inin_{\wt}(G_j).
\end{equation}

\begin{lemma}\label{lem:laplace-equivalence}
One has
\[
\alpha_0\ge m_+,
\]
and
\begin{equation}\label{eq:J0-equivalence}
J_0=0\iff\alpha_0>m_+.
\end{equation}
\end{lemma}

\begin{proof}
If \(\alpha_0<m_+\), then \(G_0\) is the unique lowest-valuation summand of \eqref{eq:P0}, impossible.  If \(\alpha_0=m_+\), the complete weight-\(m_+\) part of \eqref{eq:P0} is
\[
\inin_{\wt}(G_0)+J_0=0,
\]
so \(J_0=-\inin_{\wt}(G_0)\ne0\).  If \(\alpha_0>m_+\), the weight-\(m_+\) part consists only of the positive terms, hence \(J_0=0\).
\end{proof}

\begin{theorem}[Cancellative Leftmost]\label{thm:floor}
On the ratio-identity locus, under the hypotheses of Theorem~\ref{thm:ratio-rigidity},
\begin{equation}\label{eq:floor}
J_0=0\Longrightarrow t\in T_0^+.
\end{equation}
In fact the stronger conclusions
\begin{equation}\label{eq:floor-strong}
J_0=0,\qquad T_0^+=\{d_1,d_2\},\qquad t=d_1
\end{equation}
hold automatically.
\end{theorem}

\begin{proof}
Apply Theorem~\ref{thm:ratio-rigidity} with \(\omega(k)=k^2\).  Its
\eqref{eq:universal-floor-set} is \(T_0^+=\{d_1,d_2\}\), and its
\eqref{eq:automatic-floor-antecedent} reads \(\alpha_0>m_+\), which gives
\(J_0=0\) by \eqref{eq:J0-equivalence}.  The leftmost positive supported
shift is \(d_1\), since \eqref{eq:support-prefix} makes \(G_1\ne0\) and
\(d_1\) is the smallest positive shift by Lemma~\ref{lem:grid-order}.
\end{proof}

\subsection{Reversal pairs on the locus}\label{subsec:locus-family}

The locus is populated well past the two-support configuration, so
Theorem~\ref{thm:ratio-rigidity} and its consequences are not vacuous.  The
witnesses come in pairs, each a diagram and its row-reversal on a staircase
inner boundary, and there is one for every admissible word.

Fix \(N\ge2\), put \(\mu=(N-1,\ldots,1,0)\), and for a word
\(w=(w_0,\ldots,w_{N-1})\) with entries in \(\{2,3\}\) write
\(X_w=(w+\mu)/\mu\), the sum taken coordinatewise.  Each \(X_w\) is a connected
skew diagram with row lengths \(w\), since its adjacent overlaps are
\(w_{i+1}-1\ge1\).  Write \(w^{*}\) for the reversed word, and abbreviate
\[
D(\rho)=\det\bigl(h_{\rho_i+2(j-i)}\bigr),
\]
the determinant running over the length of \(\rho\).

The mechanism is a determinant identity for words in two formal series, and it
is insensitive to what the series are.  The Toeplitz-minor correspondence and
the row/column reversal are classical; see \cite{GarciaGarcia2020ToeplitzMinors}.
What is proved here is the binary-word identity \eqref{eq:word-minor}, uniform
in the word, and its consequence for the reversal pairs below.

\begin{lemma}\label{lem:two-series-reversal}
Let \(\mathcal A=\sum_{i\ge0}\alpha_iz^i\) with \(\alpha_0=1\) and
\(\mathcal B=\sum_{i\ge0}\beta_iz^i\) be formal power series over a commutative
ring, and let \(\mathcal Q=\mathcal B/\mathcal A=\sum_{i\ge0}q_iz^i\), with
\(q_k=0\) for \(k<0\).  For a word \(x=(x_1,\ldots,x_\ell)\) with entries in
\(\{\mathcal A,\mathcal B\}\) let \(M_x\)
be the \((\ell+2)\times(\ell+2)\) matrix on columns \(0,\ldots,\ell+1\) whose first
two rows are the coefficient sequences of \(\mathcal B\) and \(\mathcal A\) and whose
\((i+2)\)nd row is that of \(x_i\) shifted right by \(i\).  With
\(J=\{i:x_i=\mathcal B\}\),
\begin{equation}\label{eq:word-minor}
\det M_x=(-1)^{\ell+1-|J|}\det\bigl(q_{c-r}\bigr)_{r\in\{0\}\cup J,\ c\in J\cup\{\ell+1\}},
\end{equation}
and consequently \(\det M_x=\det M_{x^{*}}\).
\end{lemma}

\begin{proof}
Right multiplication by the inverse of the upper triangular Toeplitz matrix of
\(\mathcal A\) is unimodular, since \(\alpha_0=1\).  It carries each
\(\mathcal A\)-row to a coordinate row and each \(\mathcal B\)-row to the
corresponding shift of \(\mathcal Q\).  The coordinate row \(e_k\) sits in row
\(k+1\), so deleting the coordinate rows from the largest \(k\) down leaves the
positions of the others unchanged and each contributes \((-1)^{(k+1)+k}=-1\).
There are \(\ell+1-|J|\) of them, which is the sign in \eqref{eq:word-minor};
the rows that survive are those carrying the shifts \(\{0\}\cup J\) and the
columns that survive are \(J\cup\{\ell+1\}\).

Reversal replaces \(J\) by \(\{\ell+1-i:i\in J\}\), and \(|J|\) is unchanged, so
the sign is unchanged.  Substituting \(r\mapsto\ell+1-c\) and \(c\mapsto\ell+1-r\)
leaves \(c-r\) fixed and exchanges the two index sets, so the reversed minor is
the transpose of the original with both orders reversed.  Transposition preserves
the determinant, and reversing the order of the rows and of the columns
contributes the same sign twice.
\end{proof}

\begin{proposition}\label{prop:staircase-reversal}
Let \(w\in\{2,3\}^N\) with \(w_0=w_{N-1}=2\).  Then \(X_w\) and \(X_{w^{*}}\) have a
common inner boundary and a common cap grid, and satisfy the ratio identity
\eqref{eq:ratio-identity} at \((r,a)=(2,1)\).  If \(w\ne w^{*}\) the two diagrams
are distinct and \(s_{X_w}\ne s_{X_{w^{*}}}\).
\end{proposition}

\begin{proof}
Both words begin with \(2\), so both parents have \(L=N+1\), and by
\eqref{eq:dexplicit} both cap grids are \(d_j=2j\).  By \eqref{eq:Adef} the row
parameters are \(A_i=2i+2-w_i\), so \eqref{eq:capmatrix} reads
\(M_{ij}=h_{w_i+2(j-i-1)}\).  Adjoining the row \((h_{r+d_j})_j\) on top and
expanding along it gives
\[
P_r(C(X_w))=\sum_j C_j(X_w)h_{r+d_j}=D(r,w),
\]
the expansion being \eqref{eq:Pk} and \eqref{eq:cofactor} read backwards.  So with
\[
F(w)=h_1P_2(C(X_w))-P_3(C(X_w))
\]
one has \(F(w)=h_1D(2,w)-D(3,w)=D(1,2,w)\), the last step being the same
expansion along the adjoined row.

Now set \(\alpha_i=h_{2i}\) and
\(\beta_i=h_{2i+1}\).  In \(D(\rho)\) a row with \(\rho_i=1\) has entries
\(h_{1+2(j-i)}=\beta_{j-i}\), a row with \(\rho_i=2\) has
\(h_{2+2(j-i)}=\alpha_{j-i+1}\), and one with \(\rho_i=3\) has
\(\beta_{j-i+1}\); that is \(\mathcal B\) shifted by \(i\) in the first case and
\(\mathcal A\) or \(\mathcal B\) shifted by \(i-1\) in the other two.  For
\(\rho=(1,2,w)\) the two prefixed letters therefore give the rows \(\mathcal B\)
and \(\mathcal A\) at shift \(0\), and \(w_t\) gives its letter at shift \(t+1\);
so \(D(1,2,w)=\det M_w\) with \(\ell=N\).
Lemma~\ref{lem:two-series-reversal} gives \(F(w)=F(w^{*})\), and
\(w_{N-1}=2\) is what allows the same identification for \(w^{*}\).  Since
\(P_k\) is linear in its argument, \(F(w)-F(w^{*})=h_1P_2(G)-P_3(G)\) for
\(G=C(X_w)-C(X_{w^{*}})\), which is the ratio identity at \((r,a)=(2,1)\).

If \(w\ne w^{*}\) the two outer boundaries differ, so the diagrams are distinct.
Were they skew-equivalent, Corollary~\ref{cor:boundary-separation} would forbid
the ratio identity just established.
\end{proof}

The number of such pairs grows with \(N\): the interior carries \(2^{N-2}\) words,
of which \(2^{\lceil(N-2)/2\rceil}\) are palindromic, so there are
\(\bigl(2^{N-2}-2^{\lceil(N-2)/2\rceil}\bigr)/2\) unordered pairs at length \(N\),
every one of them a distinct non-equivalent pair on the locus.

\begin{corollary}\label{cor:locus-family}
For every \(k\ge1\), with \(N=2k+2\) and row-length words
\[
\rho(A_k)=(2,2^k,3^k,2),
\qquad
\rho(B_k)=(2,3^k,2^k,2)
\]
added coordinatewise to \(\mu=(N-1,\ldots,1,0)\), every clause of
Theorem~\ref{thm:ratio-rigidity} holds at \((r,a)=(2,1)\), with
\(\supp G=\{0,2,4,\ldots,4k\}\), so \(m=2k\) positive supported
shifts, and \(T_0^+=\{d_1,d_2\}\), \(\alpha_0=13k+20\), \(\alpha_0-m_+=8\).
\end{corollary}

\begin{proof}
The two words are reversals of one another and carry \(2\) at both ends, so
Proposition~\ref{prop:staircase-reversal} applies.  For the support, note that
\(M_{ij}=h_{\rho_i+2(j-i-1)}\) vanishes for \(j<i\), since \(\rho_i\le3\); deleting
column \(j\ge1\) therefore splits the determinant and factors the cofactor as
\(C_j(\rho)=(-1)^j\bigl(\prod_{i=1}^{j-1}h_{\rho_i-2}\bigr)D(\rho_j,\ldots,\rho_{N-1})\).
For \(2\le j\le2k\) the two parents' diagonal monomials differ, with square
weights separated by \(4\min(j-1,2k+1-j)>0\), so those coordinates do not cancel;
the last two coordinates agree, being \(-h_1^kh_2\) and \(h_1^k\) on both sides.
Theorem~\ref{thm:ratio-rigidity} supplies \(G_0,G_1\ne0\).  The stated
\(\alpha_0\) follows from \(\wt(G_2)=13k-4\), and the gap is
\eqref{eq:square-gap} at \(\gamma=r=2\) and \(d_1=2\).
\end{proof}

The members of each pair are \emph{not} skew-equivalent: the ratio identity holds
while \(s_{A_k}\ne s_{B_k}\).  So the locus contains pairs with \(2k\) positive
supported shifts for every \(k\ge1\), while nothing here concerns equal pairs.  The locus is a
condition on the capped functions, not on the parents.

\begin{remark}\label{rem:staircase-boundary}
The boundary letter is a hypothesis and not a normalization.  Let
\(w\in\{2,3\}^N\) be nonpalindromic with \(w_0=w_{N-1}=3\), and put
\(G=C(X_w)-C(X_{w^{*}})\).  The grid is then \(d_0=0\) and \(d_j=2j+1\) for
\(1\le j\le N\), so the identification \(P_r(C(X_w))=D(r,w)\) is unavailable.
In its place, first-column expansion gives
\[
h_1P_r(C(X_w))=D(r+1,w)+(h_1h_r-h_{r+1})D(w).
\]
Lemma~\ref{lem:two-series-reversal} still gives \(D(1,2,w)=D(1,2,w^{*})\).
Expanding first columns and using the preceding formula at \(r=1,2\) yields
\[
h_1P_2(G)-h_1^2P_1(G)+e_3G_0=0,
\qquad
e_3=h_1^3-2h_1h_2+h_3 .
\]
If \(G_0=0\), canceling \(h_1\) gives \(P_2(G)=h_1P_1(G)\), contrary to
Corollary~\ref{cor:boundary-separation} for these distinct Schur-equal parents.
So \(G_0\ne0\) and the parents are inequivalent.

Suppose now that \(P_3(G)=h_1P_2(G)\).  Theorem~\ref{thm:ratio-rigidity} and
\eqref{eq:square-gap} give \(\alpha_1=\alpha_2=u=m_+\) and \(\alpha_0=u+12\).
In \(P_1(G)\) column one uniquely minimizes, at weight \(u+7\), against
\(u+13\) at column zero and at least \(u+11\) later; every summand of
\(P_2(G)\) has weight at least \(u+16\); and \(\wt(e_3)=3\).  The three terms
of the displayed relation then have weights at least \(u+17\), exactly
\(u+9\), and exactly \(u+15\), so the relation has a unique minimum and cannot
vanish.  No nonpalindromic word on this boundary therefore satisfies the ratio
identity at \((r,a)=(2,1)\).  The smallest witness is \(w=(3,2,3,3)\), where
the square-weight initial form of \(h_1P_2(G)-P_3(G)\) is
\(-h_1^3h_2h_4h_5\).
\end{remark}

\section{Two-cap exclusion}\label{sec:two-cap}

The preceding rigidity supplies the valuation control for two-cap exclusion.

\subsection{Two caps cannot coexist}

The cap relations available to a cofactor defect come in two forms.  A \emph{vanishing cap} at shift \(k\) is the relation \(P_k(G)=0\); the case \(k=0\) is automatic, being the row-\(0\) case of \eqref{eq:common-row-relation} together with \(A_0=0\).  A \emph{ratio cap} at \((r,a)\) is the identity \eqref{eq:ratio-identity}.  A vanishing cap at \(a>0\) is the case \(r=0\) of a ratio cap, since \(P_0(G)=0\) makes \(P_a(G)=h_aP_0(G)\) read \(P_a(G)=0\).

Throughout this subsection \(\omega(k)=k^2\) and
\begin{equation}\label{eq:Phi}
\Phi_k(j)=\wt(G_j)+(k+d_j)^2
=
\alpha_j+2kd_j+k^2,
\end{equation}
with \(\Phi_k(j)=+\infty\) when \(G_j=0\).  The second equality uses \(d_0=0\) and holds for every \(j\).  The quantity \(\Phi_k(j)\) is the square weight of the \(j\)-th summand of \(P_k(G)=\sum_j G_jh_{k+d_j}\), by Lemma~\ref{lem:valuation-laws}.

\begin{theorem}[Two-cap exclusion]\label{thm:two-cap}
Let \(G\) be a defect on an ordered cap grid with \(d_0=0\), let \(r\ge0\) and \(0<a<b\), and suppose
\begin{equation}\label{eq:two-caps}
P_{r+a}(G)=h_aP_r(G),
\qquad
P_{r+b}(G)=h_bP_r(G).
\end{equation}
Put \(\gamma=a\) if \(r=0\) and \(\gamma=r\) if \(r>0\).  If \(G_0\ne0\) and
\begin{equation}\label{eq:floor-input}
\alpha_0-\alpha_j\le2\gamma d_j
\qquad\text{for every }j\text{ with }G_j\ne0,
\end{equation}
then no such \(G\) exists.
\end{theorem}

\begin{proof}
Suppose first \(r=0\).  Then \(P_0(G)=0\) by \eqref{eq:P0-vanishes}, which holds because \(G\) is the defect of two connected parents and not merely a vector, so \eqref{eq:two-caps} reads \(P_a(G)=P_b(G)=0\).  For every \(j\) with \(G_j\ne0\), \(\Phi_a(j)-\Phi_a(0) = \alpha_j-\alpha_0+2ad_j \ge0\)
by \eqref{eq:floor-input}, while
\begin{equation}\label{eq:cap-shift}
\Phi_b(j)-\Phi_b(0)
=
\bigl(\Phi_a(j)-\Phi_a(0)\bigr)+2d_j(b-a),
\end{equation}
directly from \eqref{eq:Phi}.  For \(j\ne0\) one has \(d_j>0\) and \(b>a\), so the right-hand side of \eqref{eq:cap-shift} is strictly positive.  Thus the \(j=0\) summand of \(P_b(G)=\sum_jG_jh_{b+d_j}\) has strictly smallest square weight among the nonzero summands, and \(\wt(P_b(G))=\Phi_b(0)<\infty\) by Lemma~\ref{lem:valuation-laws}.  Hence \(P_b(G)\ne0\), contradicting the second cap.

Suppose now \(r>0\), and let \(k\in\{a,b\}\).  For every \(j\ne0\) with \(G_j\ne0\),
\[
\Phi_{r+k}(j)-\Phi_{r+k}(0)
=
\alpha_j-\alpha_0+2(r+k)d_j
\ge
-2rd_j+2(r+k)d_j
=
2kd_j
>0,
\]
using \eqref{eq:floor-input} with \(\gamma=r\).  So \(j=0\) is again the strict unique minimizer, and
\[
\wt\bigl(P_{r+k}(G)\bigr)
=
\Phi_{r+k}(0)
=
\wt(G_0)+(r+k)^2.
\]
This is finite, so \(P_{r+k}(G)\ne0\) and, by \eqref{eq:two-caps} and Lemma~\ref{lem:valuation-laws}, \(P_r(G)\ne0\) with
\[
\wt\bigl(P_r(G)\bigr)
=
\wt\bigl(P_{r+k}(G)\bigr)-k^2
=
\wt(G_0)+(r+k)^2-k^2
=
\wt(G_0)+r^2+2rk.
\]
The left-hand side does not depend on \(k\), so \(2ra=2rb\), and \(r>0\) forces \(a=b\).
\end{proof}

\begin{corollary}\label{cor:floor-supplies}
Let \(A\) and \(B\) be connected skew diagrams with a common inner boundary, let \(G=C(A)-C(B)\ne0\), and assume the ratio identity \eqref{eq:ratio-identity} at \((r,a)\).  Then \(G_0\ne0\) and \eqref{eq:floor-input} holds.
\end{corollary}

\begin{proof}
Corollary~\ref{cor:cap-fixing} puts the two diagrams on a common cap grid, so Theorem~\ref{thm:ratio-rigidity} applies with \(\omega(k)=k^2\).  It gives \(\mathcal M_{-1}^\omega=\{0,1\}\) in \eqref{eq:virtual-edge}; a column at which the defect vanishes has infinite tilted cost and cannot minimize, so \(G_0\ne0\).  It also gives \(\alpha_0-m_+=2\gamma d_1\) by \eqref{eq:square-gap}.  For any \(j\ne0\) with \(G_j\ne0\) we have \(\alpha_j\ge m_+\) and \(d_j\ge d_1\) by Lemma~\ref{lem:grid-order}, hence
\[
\alpha_0-\alpha_j
\le
\alpha_0-m_+
=
2\gamma d_1
\le
2\gamma d_j.
\]
For \(j=0\) the inequality reads \(0\le0\).
\end{proof}

\begin{corollary}\label{cor:double-cap-kills}
Let \(A\) and \(B\) be connected skew diagrams, with neither boundary prescribed, and let \(G\) be their cofactor defect \eqref{eq:defect} on the union of their cap grids.  If \(G\) satisfies two ratio identities \eqref{eq:two-caps} at \((r,a)\) and \((r,b)\) with \(0<a<b\), then \(A\) and \(B\) are translates; in particular \(G=0\).
\end{corollary}

\begin{proof}
Put \(W_X=X\) when \(r=0\) and \(W_X=X\cdot(r)\) when \(r>0\).  At \(r=0\) the relation \eqref{eq:P0-vanishes} turns the two identities into \(s_{W_A\cdot(a)}=s_{W_B\cdot(a)}\) and \(s_{W_A\cdot(b)}=s_{W_B\cdot(b)}\); at \(r>0\) the identity \eqref{eq:iterated-cap} gives the same two equalities.  Corollary~\ref{cor:two-positive-extensions} applied to \(W_A\) and \(W_B\) makes them translates, and deleting the attached row when \(r>0\) makes \(A\) and \(B\) translates.  Equal cofactor vectors then give \(G=0\).

With a common inner boundary the conclusion also follows from Theorem~\ref{thm:two-cap}, by Corollary~\ref{cor:floor-supplies} applied at \((r,a)\).
\end{proof}

\begin{corollary}\label{cor:unique-ratio}
Let \(A\) and \(B\) be connected skew diagrams with a common inner boundary, and let \(G=C(A)-C(B)\ne0\).  Then at most one pair \((r,a)\) with \(r\ge0\) and \(a>0\) satisfies \eqref{eq:ratio-identity}.
\end{corollary}

\begin{proof}
A ratio identity fixes the cap by Corollary~\ref{cor:cap-fixing}, and by \eqref{eq:square-gap} the number
\[
\gamma=\frac{\alpha_0-m_+}{2d_1}>0
\]
is determined by \(G\) alone, while equalling \(r\) when \(r>0\) and \(a\) when \(r=0\).  Two identities with positive base therefore share that base, and Theorem~\ref{thm:two-cap} forces their cap lengths to agree; two identities with base zero share their cap length immediately.

One identity of each kind remains to be excluded.  Their parameters would be \((0,\gamma)\) and \((\gamma,b)\) for some \(b>0\).  Since \(P_0(G)=0\) by \eqref{eq:P0-vanishes}, the first gives \(P_\gamma(G)=0\) and the second then gives \(P_{\gamma+b}(G)=h_bP_\gamma(G)=0\).  These are two ratio identities at base zero with cap lengths \(\gamma<\gamma+b\), which Theorem~\ref{thm:two-cap} excludes.
\end{proof}

\section{The seam chain and the composition problem}\label{sec:application}

The reading of the adjacent-edge chain below is an interpretation; the
reconstruction theorem is unconditional.

\subsection{The cancellation chain and the composition problem}\label{subsec:seam-interpretation}

For an actual equal pair, Theorem~\ref{thm:equal-pair-rigidity} gives the
chain of minimizing column pairs
\begin{equation}\label{eq:tropical-seam-chain}
\{1,2\},\{2,3\},\ldots,\{v,v+1\}.
\end{equation}
We call these pairs \emph{seams}.  They record leading
cancellations in the exact common-row relations.  The further pair
\(\{0,1\}\) belongs to the virtual row on the ratio-identity locus alone
and cannot occur here, because \(G_0=0\)
(Remark~\ref{rem:no-virtual-row}).

The leading geometry also cannot distinguish the two crossing pairings of
four columns, and that is an algebraic fact rather than an analogy.
Calling what it misses an orientation is the analogy.  Consider
an \(m\times(m+2)\) staircase \((h_{d_j-A_i})\) with
\(A\) and \(d\) strictly increasing and \(A_t\le d_t\) for \(0\le t<m\),
so that every increasing matching is feasible, and for \(i<j\) let
\(P_{ij}\) be the maximal minor deleting columns \(i\) and \(j\).  By
Theorem~\ref{thm:monge-minor} each has a unique initial monomial
\(M_{ij}\), of weight \(w_{ij}\).

\begin{proposition}\label{prop:caterpillar}
For \(i<j<k<l\), on such a staircase,
\[
M_{ik}M_{jl}=M_{il}M_{jk}
\]
as monomials, hence
\[
w_{ik}+w_{jl}=w_{il}+w_{jk}
<
w_{ij}+w_{kl}.
\]
\end{proposition}

\begin{proof}
For deletion of \(a<b\) the increasing matching sends row \(t\) to \(t+\mathbf1_{t\ge a}+\mathbf1_{t\ge b-1}\).  In \(P_{ik}P_{jl}\) and \(P_{il}P_{jk}\) the two determinants therefore retain the same unordered pair of columns in \emph{every} row, so the increasing matchings contribute the same multiset of subscripts row by row.  This gives the monomial identity and the equality of weights.  The comparison has to be rowwise: all three pairings retain the same \emph{total} multiset of columns, so that count alone cannot separate the crossing pairings from the separated one.  For the strict inequality, the separated pairing \(\{i,j\},\{k,l\}\) assigns to one determinant a matching that is crossed relative to the other two; uncrossing it, as in Lemma~\ref{lem:strict-monge}, strictly decreases the weight.
\end{proof}

Three pairings of four columns into two pairs exist, and the initial-form
geometry sees only two of them: it strictly rejects the separated pairing
and cannot distinguish the two crossing ones.  We read that
indistinguishability as an orientation datum invisible to the leading
geometry.

That is consistent with factor reversal being the residual freedom in the
McNamara--van Willigenburg formulation, and it is not a proof of it.
Nothing here identifies an adjacent pair of minimizing columns with a legal
marked overlap, produces a smaller equality to descend to, or exhibits an
independently reversible diagrammatic factor; a four-point identity for
initial minors supplies none of those data.  What is proved instead is
boundary reconstruction.  For parents sharing an inner boundary, separation
follows from the nonvanishing of the first positive defect coordinate; for
arbitrary boundaries it follows from the prefix-minor induction of
Lemma~\ref{lem:two-cofactor-reconstruction}.  Section~\ref{sec:scope} states what a
classification argument would still require.

\subsection{Composition and reversal}

Two operations generate the equalities the conjecture predicts.  The first is reversal.

\begin{definition}[Factor reversal]\label{def:reversal}
For a skew diagram \(F\), write \(F^{*}\) for its rotation by \(180\) degrees.  A skew diagram and its rotation are always skew-equivalent, and we call the passage from \(F\) to \(F^{*}\) \emph{reversal}.
\end{definition}

The second is composition, which we recall in the form of \cite[Defs.~3.5, 3.11]{McNamara2009Classification}.  A skew diagram \(W\) \emph{lies in the top} (respectively \emph{bottom}) of a skew diagram \(E\) if \(W\) occurs as a connected subdiagram of \(E\) containing the northeasternmost (respectively southwesternmost) cell of \(E\); the empty diagram lies in both.  If \(W\) lies in the top of \(E_1\) and in the bottom of \(E_2\), the \emph{amalgamation} \(E_1\odot_W E_2\) is obtained from the disjoint union of \(E_1\) and \(E_2\) by identifying the two copies of \(W\).  When \(W\) lies in both the top and the bottom of \(E\), write \(E=W\!\cdot\!O\!\cdot\!W\) for the decomposition in which \(O\) is what remains after deleting both copies.  Four conditions are imposed on the pair \((E,W)\).  Hypotheses I--III are \cite[Hyp.~3.8]{McNamara2009Classification}: \(W\) is maximal among diagrams occupying its diagonals that lie in both the top and the bottom of \(E\); the two copies of \(W\) are separated by at least one diagonal; and the complement in \(E\) of either copy is a connected skew diagram.  Hypothesis IV is the separate condition \cite[Hyp.~3.22]{McNamara2009Classification}, a nonadjacency requirement imposed on the infinite amalgamation \(E\!\odot\!W^{\infty}\), and stated there for auxiliary diagrams derived from \(O\) and \(W\) rather than for \(O\) and \(W\) themselves.

\begin{definition}[MvW composition]\label{def:composition}
For skew diagrams \(D\) and \(E=W\!\cdot\!O\!\cdot\!W\) subject to those hypotheses, the composition \(D\circ_W E\) is the skew diagram assembled from one copy \(E_d\) of \(E\) per cell \(d\) of \(D\), adjacent copies being amalgamated along \(W\) or translated one position according to how \(O\) meets the two copies of \(W\) in \(E\).  That splits into four cases, and the assembly rule for each is \cite[Def.~3.11]{McNamara2009Classification}.
\end{definition}

\begin{definition}[Generated by composition and reversal]\label{def:generated}
An equality \(s_D=s_E\) of skew Schur functions is \emph{generated by McNamara--van Willigenburg composition and factor reversal} if there are \(r\ge1\) and skew diagrams \(E_1,\ldots,E_r\), with \(E_i=W_i\!\cdot\!O_i\!\cdot\!W_i\) subject to the hypotheses above for \(i\ge2\), such that
\[
D=\bigl(\cdots((E_1\circ_{W_2}E_2)\circ_{W_3}E_3)\cdots\bigr)\circ_{W_r}E_r,
\qquad
E=\bigl(\cdots((E_1'\circ_{W_2'}E_2')\circ_{W_3'}E_3')\cdots\bigr)\circ_{W_r'}E_r',
\]
where \(E_1'\) is \(E_1\) or \(E_1^{*}\) and, for \(i\ge2\), the pair \((E_i',W_i')\) is \((E_i,W_i)\) or \((E_i^{*},W_i^{*})\).
\end{definition}

If each \(E_i\) additionally satisfies Hypothesis V of \cite[Hyp.~3.27]{McNamara2009Classification} -- one copy of \(W\) has a single cell adjacent to \(O\) -- then every such pair is skew-equivalent by \cite[Thm.~5.3]{McNamara2009Classification}.  Hypothesis V is not needed for that direction.  Olson-Harris removes it by Hopf-algebraic methods in place of the determinantal route, building on \cite{Yeats2017HopfSchur}: \cite[Thm.~5.3.9]{OlsonHarris2024HopfApplications} gives \(D\circ_WE\sim D'\circ_WE\) whenever \(D\sim D'\) and \((E,W)\) satisfies Hypotheses I--IV, and the tower form follows there.  The composition identity of \cite[Conj.~3.26]{McNamara2009Classification} is itself \cite[Thm.~5.6.4]{OlsonHarris2024HopfApplications}, proved there without Hypothesis V; the necessity clause that conjecture attaches to Hypothesis IV is a separate assertion, not settled by either theorem.  The converse, that every equality is generated in this way, is \cite[Conj.~5.7]{McNamara2009Classification}; nothing below asserts it.

\begin{remark}\label{rem:factorwise}
The freedom in Definition~\ref{def:generated} is genuinely factor-wise.  On ribbons, where the classification is the theorem of \cite{Billera2006DecomposableCompositions}, \((1,2)\circ(1,2)=(1,2,1,3,2)\) is equivalent to \((2,1)\circ(1,2)=(1,3,2,1,2)\), which is neither it nor its reversal \((2,3,1,2,1)\): reversing the two factors independently gives a class of four ribbons, whereas global reversal gives classes of size at most two.  Over all compositions of size at most nine this happens only at nine, by exhaustive comparison of the ribbon Schur expansions at each size.  The orientation freedom left by \eqref{eq:tropical-seam-chain} therefore cannot be tightened.
\end{remark}

\subsection{Every positive cap separates an equal pair}\label{subsec:cap-separation}

\begin{corollary}\label{cor:boundary-separation}
In the situation of Theorem~\ref{thm:equal-pair-rigidity}, for every integer \(a>0\),
\[
P_a(G)\ne0 .
\]
More precisely, with \(\varepsilon_1M_1\) as in Proposition~\ref{prop:monomial-prefix},
\begin{equation}\label{eq:cap-leading-monomial}
\inin_\omega P_a(G)=\varepsilon_1h_{a+d_1}M_1
\end{equation}
for every admissible convex weight, and for \(r>0\)
\begin{equation}\label{eq:ratio-leading-monomial}
\inin_\omega\bigl(P_{r+a}(G)-h_aP_r(G)\bigr)
=-\varepsilon_1h_ah_{r+d_1}M_1 .
\end{equation}
In particular no ratio identity at \(r\ge0\), \(a>0\) holds on a distinct
Schur-equal pair with a common inner boundary.
\end{corollary}

\begin{proof}
For each positive supported coordinate, Lemma~\ref{lem:valuation-laws} gives
\[
\wto\bigl(G_jh_{a+d_j}\bigr)
=
\alpha_j^\omega+\omega(a+d_j)-\omega(d_j),
\]
and the increment \(d\mapsto\omega(a+d)-\omega(d)\) is strictly increasing because \(\Delta\omega\) is.  By \eqref{eq:prefix-profile} the coordinates \(j\le v+1\) have \(\alpha_j^\omega\ge\alpha_1^\omega=m_+^\omega\), and by the closing statement of Theorem~\ref{thm:cap-fixed-rigidity} every later positive coordinate has normalized action strictly larger than \(\alpha_{v+1}^\omega\); so \(\alpha_j^\omega\ge m_+^\omega\) at every positive coordinate, while \(G_0=0\) removes the zeroth summand altogether.  Column one is therefore the unique minimizing summand, which is \eqref{eq:cap-leading-monomial}.

For \(r>0\), positive strictly increasing increments give \(\omega(x+y)>\omega(x)+\omega(y)\) for \(x,y\ge1\), so
\[
\inin_\omega\bigl(h_{r+a+d_j}-h_ah_{r+d_j}\bigr)=-h_ah_{r+d_j},
\]
and the same comparison, now with the shift \(r\), again isolates column one.  This is \eqref{eq:ratio-leading-monomial}.  At \(r=0\) the ratio difference is \(P_a(G)\) itself, since \(P_0(G)=0\) by \eqref{eq:P0-vanishes}.
\end{proof}

By \eqref{eq:Pk-is-capped} the first assertion says that attaching a row of
any positive length in the northeast position distinguishes two equal
parents.  Equations \eqref{eq:cap-leading-monomial} and
\eqref{eq:ratio-leading-monomial} name the one surviving initial monomial in
each boundary response; they do not distinguish the original functions,
which are equal.

\section{Scope and remaining questions}\label{sec:scope}

The examples below distinguish cofactor realizability from formal cap
relations, and valuation-specific collisions from the boundary hypotheses.
The majorization results and their moving-cap limitation are collected in
Appendix~\ref{app:structural}.

\subsection{The ratio identity alone does not give Cancellative Leftmost}\label{subsec:elbow}

Theorem~\ref{thm:floor} is stated on the ratio-identity locus, and it is fair to ask whether the two relations that locus supplies are the whole hypothesis: the vanishing \eqref{eq:P0-vanishes} and the ratio identity, read as conditions on a formal vector of \(\Lam^3\) rather than on the cofactor defect of a pair of parents.  They are not, and the obstruction is explicit.

\begin{proposition}\label{prop:elbow}
Fix \(x<y<z\) and a cap \((r,a)\), write
\[
Y_d=h_{r+a+d}-h_ah_{r+d},
\]
and let \(G\) be supported on \(\{x,y,z\}\) with
\[
(G_x,G_y,G_z)
=
(h_x,h_y,h_z)\times(Y_x,Y_y,Y_z),
\]
the cross product in \(\Lam^3\).  Then
\[
P_0(G)=0
\qquad\text{and}\qquad
P_{r+a}(G)=h_aP_r(G).
\]
\end{proposition}

\begin{proof}
A cross product is orthogonal to each of its factors.  Orthogonality to \((h_x,h_y,h_z)\) is \(\sum_dh_dG_d=0\), which is \(P_0(G)=0\); orthogonality to \((Y_x,Y_y,Y_z)\) is \(\sum_dY_dG_d=0\), which is the ratio identity.
\end{proof}

When \(x>0\), these vectors satisfy both relations that the argument of Theorem~\ref{thm:floor} has available and fail its conclusion.  At \((x,y,z,r,a)=(1,2,3,1,1)\) the square-weight data are \((\alpha_x,\alpha_y,\alpha_z)=(20,18,18)\),
so \(T_0^+=\{2,3\}\) and \(J_0=0\), while the leftmost positive supported shift is \(t=1\), which is not in \(T_0^+\).  For every \(0\le x<y<z\) one has \(\alpha_y=\alpha_z\) and \(\alpha_x-\alpha_y=2\gamma(y-x)\)
with \(\gamma\) as in \eqref{eq:gamma}, so the failure persists at every \(x>0\).  At \(x=0\) it does not: the leftmost positive supported shift is then \(y\), and \(T_0^+=\{y,z\}\), so \eqref{eq:floor} holds.  The counterexample needs the zeroth shift unsupported.

That last equation is the square-gap formula \eqref{eq:square-gap} of Theorem~\ref{thm:ratio-rigidity}, with \(x\) in the role of the zeroth coordinate and \(y-x\) in the role of \(d_1\).  So the cone does not merely evade the conclusion; it reproduces the numerology the theorem predicts, which shows that \eqref{eq:square-gap} follows from the tie structure alone and carries no realizability with it.  The square-gap identity alone therefore does not certify cofactor realizability.

The converse deficiency is also real: there are Jacobi--Trudi realizable configurations with \(\lvert T_+\rvert=2\) satisfying no ratio identity at any \((r,a)\).  So neither hypothesis alone suffices, and what Theorem~\ref{thm:floor} uses beyond the ratio identity is that its input is the cofactor defect of two parents: besides \eqref{eq:P0-vanishes} its proof consumes the synchronized parent profiles, the first-split anchor, and the common-prefix kernel relations.  Remark~\ref{rem:realizability-terminal} is the same phenomenon at Theorem~\ref{thm:terminal-gap}.

\subsection{Weight collisions and coefficient cancellation}\label{subsec:weight-scope}

A supported defect coordinate can have the form \(G=-h_3^2+h_4h_1^2\),
whose image under \(h_j\mapsto q^{j^2}\) is zero: the distinct monomials
both have square weight \(18\).  This failure concerns the square-weight
collapse, not every admissible convex weight.

For an integer \(B\ge3\) put
\begin{equation}\label{eq:radix-weight}
\omega_B(0)=0,
\qquad
\omega_B(k)=B^k
\quad(k\ge1).
\end{equation}
The increments are \(\Delta\omega_B(1)=B\) and \(\Delta\omega_B(k)=B^k-B^{k-1}\) for \(k\ge2\), all positive and strictly increasing, so \(\omega_B\) is admissible and every result stated for an arbitrary admissible convex weight applies to it.  The bound \(B\ge3\) is sharp: at \(B=2\) the first two increments are both \(2\) and strict convexity fails.

Now let \(X\) have \(n\) rows.  Every monomial of a maximal minor of \eqref{eq:capmatrix} is a product of at most \(n\) generators, so each exponent is at most \(n\); if \(B>n\) then
\[
\nu_{\omega_B}\Bigl(\prod_j h_j^{e_j}\Bigr)
=
\sum_j e_jB^j
\]
is the base-\(B\) numeral with digits \(e_j\), and no carrying occurs.  Hence \(\omega_B\) is \emph{injective} on exactly the monomials these determinants produce: distinct monomials never share a weight.  The radix weight therefore separates the distinct monomials that the square weight can conflate.

For the coordinate above, \(B=7\) gives weights \(686\) and \(2415\).
Higher cap-fixed strata can likewise contain incomparable monomials of
equal square weight, as Remark~\ref{rem:higher-strata} exhibits; radix
weights separate those collisions.  They do not restore absent terms.
The minors \(P_{ij}\) of Proposition~\ref{prop:caterpillar} satisfy the
three-term Pl\"ucker relation \(P_{ij}P_{kl}-P_{ik}P_{jl}+P_{il}P_{jk}=0\)
for \(i<j<k<l\), and the three products can be pairwise distinct.  The
present argument requires neither coordinatewise square-collapse
faithfulness nor a classification of those relations: it uses exact
relations to locate their minima.

\subsection{The strict floor condition is essential}\label{subsec:floor-margin}

The conclusion about the leftmost positive coordinate fails as soon as the strict floor condition is replaced by equality, and it fails inside the class rather than outside it: with a common inner boundary, a common cap grid, equal row-length multisets and only three rows.  Take
\[
A=(4,4,2)/(1,1,0),
\qquad
B=(4,3,3)/(1,1,0).
\]
Both parents are connected, both have row-length multiset \(\{3,3,2\}\), and they share the cap grid \((0,3,4,6)\).  Their positive defect coordinates are
\[
G_1=h_5-h_1h_4,
\qquad
G_2=h_2^2-h_1h_3+h_4,
\qquad
G_3=h_1^2-2h_2,
\]
with \(G_0=-h_3G_1-h_4G_2-h_6G_3\).  For square weight, \((\alpha_0,\alpha_1,\alpha_2,\alpha_3)=(24,26,24,38)\),
so \(\alpha_0=m_+=24\), while the leftmost positive supported shift is \(3\) and \(T_0^+=\{4\}\).

The discrepancy is unbounded with the number of rows and the number of supported coordinates held fixed.  For an integer \(k\ge1\) put
\[
\mu^{[k]}=(3k-2,2k-1,0),
\qquad
A_k=(6k-2,5k-1,2k)/\mu^{[k]},
\qquad
B_k=(6k-2,4k-1,3k)/\mu^{[k]},
\]
with cap grid \((0,3k,4k,6k)\) and row-parameter sequences \((0,k,4k)\) and \((0,2k,3k)\).  Every Jacobi--Trudi subscript is \(k\) times the corresponding subscript above, so the defect is obtained by the substitution \(h_j\mapsto h_{kj}\) and \((\alpha_0,\alpha_1,\alpha_2,\alpha_3)=k^2(24,26,24,38)\).
The leftmost positive action exceeds the positive floor by \(2k^2\).  No repaired statement carrying a fixed slack term can therefore exist, and the dilation proves this for every \(k\) rather than extrapolating from a sweep.

\subsection{The fingerprint does not identify an equality}\label{subsec:fingerprint-not-classification}

The cap-fixed antecedent does not imply equality, even for a legal unit transfer whose row and column overlap partitions agree on both sides.  For
\[
X=(7,5,4)/(4,2,0),
\qquad
Y=(7,6,3)/(4,2,0),
\]
the cap grid is \((0,3,6,9)\) and the square-weight defect data are
\[
\alpha_0=52,
\qquad
\alpha_1=\alpha_2=46,
\qquad
\alpha_3=+\infty,
\]
so Theorem~\ref{thm:cap-fixed-rigidity} applies; yet \(s_X-s_Y=h_3h_7-h_4h_6\ne0\).
The adjacent-edge condition is a consequence of equality in Theorem~\ref{thm:equal-pair-rigidity} and not a substitute for it, so no argument may stop at the fingerprint and conclude that two diagrams are skew-equivalent.

\subsection{Toward the general MvW conjecture}\label{subsec:future}

For general equal skew Schur functions both boundaries may change, so a
global argument first needs comparable coordinates.  The moves fall
into three pairs: normalize the boundaries onto a common ordered
Strict-Monge grid and locate their first divergence; establish
synchronization, first-split localization and no-backflow, then extract an
adjacent seam from the exact relations; identify that seam with an MvW
composition boundary along which equality descends, then recurse, with
reversal accounting for the residual orientations.

The recursion needs a separate base case: a composition-irreducible skew
diagram must be determined up to rotation by its skew Schur function.
Seam extraction does not provide this when no further cut exists.  Thus a
global proof requires both a universal seam-descent mechanism and an
irreducible classification; neither is established here.

One route is closed at the outset, and it is the one the cap relations
invite.  By Corollary~\ref{cor:boundary-separation} a distinct equal pair
with a common inner boundary satisfies no ratio identity at any \(r\ge0\),
\(a>0\); the ratio-identity locus of Section~\ref{sec:full-rigidity} is
therefore disjoint from actual equality, and
Subsection~\ref{subsec:locus-family} exhibits inequivalent pairs on it with
\(2k\) positive supported shifts for every \(k\ge1\).  So a reduction of an equality to the
double-cap configuration cannot take the equal pair itself, however its cap
data is normalized.  It has to produce different diagrams and keep a
relationship to the original pair, and that construction is what a
classification proof still owes.  Boundary determination does not supply
it: Theorem~\ref{thm:boundary-reconstruction} recovers a diagram from two
skew Schur functions, and a factorization has to be recovered from one.

\section{Concluding remarks}\label{sec:conclusion}

Two boundary observations recover both boundaries of a connected skew diagram
up to translation, prefix-minor eliminations returning the row and column
parameters in turn.  Skew Schur equality with a common inner boundary forces
in addition its own cap normalization and its own adjacent-edge fingerprint,
and the leading forms along that chain are single monomials that do not move
with the valuation.  The ratio
identity supplies the same chain by a different route, with a virtual row
in place of the vanishing zeroth coordinate, and two such identities
exclude a nonzero defect outright.  Two extensions of distinct positive
lengths recover the diagram as well, without the original function and
without the lengths
(Corollary~\ref{cor:two-positive-extensions}).

The mechanism controls where leading cancellation occurs, not which
coefficients survive.  The first-split barrier makes this control local:
later changes cannot enter the leading common-prefix equations.  The
resulting serial seam geometry is the shape McNamara--van Willigenburg
composition has, and it is rigid without the higher cancellation algebra
being simple.  What it does not supply is a factorization: the leading
geometry cannot distinguish the two crossing pairings of four columns,
which we read as the orientation freedom factor reversal names.

\appendix

\section{Structural consequences of the Strict-Monge geometry}\label{app:structural}

None of the statements below is used in the proof of Theorem~\ref{thm:boundary-reconstruction}.  Each is a consequence of the same Strict-Monge geometry, recorded on its own.  Theorem~\ref{thm:abstract-monge} holds for an arbitrary strictly Monge cost on a Ferrers support, with no symmetric-function input, so the skew Schur setting is an instance of it rather than a hypothesis in it.

\subsection{Majorization of the peel multisets, and its limits}\label{subsec:majorization}

For a parent \(X\) and a column \(j\), the \emph{peel multiset} is
\[
\bigl\{d_{k_i(j)}-A_i^X:0\le i<n\bigr\},
\]
counted with multiplicity, where \(k_i(j)\) is given by \eqref{eq:matching-index}; its positive entries are the generator subscripts of \eqref{eq:cofactor-initial}, and zero entries are retained.

For the theorem below, let \(A\) and \(B\) have a common cap grid and equal row-length multisets, let \(v\) be the first changed row, and let \(q=v+1\).  Write \(\pi_A\) and \(\pi_B\) for their peel multisets at column \(q\), arranged in decreasing order.  They have equal length and equal total sum, because the two parent cofactors have the same homogeneous degree.

\begin{theorem}\label{thm:first-split-majorization}
If \(A_v^A>A_v^B\), then
\[
\pi_A\succ\pi_B
\]
in the dominance order; if \(A_v^A<A_v^B\), the reverse dominance holds.
\end{theorem}

\begin{proof}
Assume \(A_v^A>A_v^B\); the other case is the same with the names exchanged.  By \eqref{eq:matching-index} the peel multiset at column \(0\) is the row-length multiset, and passing from column \(j\) to column \(j+1\) changes one entry only: row \(j\) moves from \(d_{j+1}-A_j\) to \(d_j-A_j\), a drop of \(d_{j+1}-d_j\).

For \(i<v\) the two parents share \(A_i\) and the grid, so those entries agree one by one; the row-length multisets agree by hypothesis.  Hence the two peel multisets at column \(v\) coincide, say in \(C\).  Row \(v\) contributes
\[
x=d_{v+1}-A_v^A
\quad\text{to \(A\)},
\qquad
y=d_{v+1}-A_v^B
\quad\text{to \(B\)},
\]
so \(C\) contains both; let \(W\) be \(C\) with one copy of each removed.  Put \(\Delta=d_{v+1}-d_v>0\).  Passing to column \(q\) lowers row \(v\) by \(\Delta\) in each parent, so
\[
\pi_A=W\uplus\{y,\,x-\Delta\},
\qquad
\pi_B=W\uplus\{x,\,y-\Delta\},
\]
where \(y>x\) because \(A_v^A>A_v^B\), and \(x-\Delta=d_v-A_v^A\ge0\) by connectedness.  Any rows changed after \(v\) are carried inside \(W\).

The two displayed pairs have equal sums, and \(\max\{x,y-\Delta\}<y=\max\{y,x-\Delta\}\), so the first majorizes the second; adjoining the common \(W\) to both preserves that.  Hence \(\pi_A\succ\pi_B\).
\end{proof}

\begin{remark}\label{rem:higher-strata}
The theorem is deliberately localized at the first split.  Fix an admissible \(\omega\) and grade the monomials of a coordinate \(G_j\) by their \(\omega\)-weight: the lowest occupied level is \(\inin_\omega G_j\), and the levels above it are the \emph{residual strata} of that coordinate.  Dominance need not survive into them, and it already fails at a fixed cap grid.  For
\[
\mu=(12,10,6,6,4,0),
\qquad
\lambda=(15,15,14,10,10,5),
\qquad
\kappa=(15,15,14,11,8,6),
\]
both parents are connected and have equal row-length multisets, and by \eqref{eq:dexplicit} they share the cap grid \((0,3,6,11,12,15,20)\), the first changed row being \(v=3\).  At \(j=6\), beyond the first split, the square-weight \(43\) stratum of \(G_6\) is \(h_1^2h_4h_5-h_3^2h_5\), whose index multisets \((5,4,1,1)\) and \((5,3,3)\), padded with zeros to a common length as above, have equal sum \(11\) and are incomparable in dominance.  Majorization here is a first-split rigidity statement, not a global anti-cancellation statement, and the moving-cap witness below bounds it.
\end{remark}

Say that \(A\) and \(B\) differ by a \emph{single non-top transfer} if their outer boundaries agree except at two rows \(r,s\ge1\), where \(\lambda_r\) is replaced by \(\lambda_r+\theta\) and \(\lambda_s\) by \(\lambda_s-\theta\) for some \(\theta>0\), both results being connected skew diagrams on the common inner boundary.  The next proposition gives comparability for such a pair; the moving-cap witness that contrasts with it follows.

\begin{proposition}\label{prop:cap-fixed-comparability}
Let \(A\) and \(B\) have a common inner boundary and differ by a single non-top transfer, so that \(\lambda_0\) is unchanged.  Then at every cap-grid shift \(d_j\) the two minimal peel multisets are comparable in dominance.
\end{proposition}

\begin{proof}
Because \(\lambda_0\) is unchanged and the inner boundary is common, the shift lists of the two parents coincide, so a common value of \(d\) deletes the same column on both sides.  By \eqref{eq:Qj} the minimal monomial at that column is indexed by the peel, and the transfer alters exactly two of its entries, replacing \(x\) and \(y\) by \(x+\theta\) and \(y-\theta\); every other entry is common.  Two multisets of equal cardinality and equal sum differing in two coordinates in this way are comparable, and adjoining a common multiset to both preserves comparability.
\end{proof}

There are cap-moving class pairs for which the two minimal peel multisets at a common shift are \((3,3)\) and \((4,1,1)\), which have equal total degree and equal square weight but are incomparable in dominance.  This does not contradict Theorem~\ref{thm:first-split-majorization}: the latter requires a common cap grid and is localized at the first split.  The witness is
\[
A=(6,6,4,1)/(2,0,0,0),
\qquad
B=(8,4,4,1)/(2,0,0,0),
\]
of fifteen cells in four rows, at the shift where
\[
C_9(A)=h_3^2-h_6,
\qquad
C_9(B)=h_4h_1^2-2h_5h_1+h_6 .
\]

\subsection{Structural byproducts}\label{subsec:byproducts}

Two consequences stated separately above are restatements and are recorded rather than numbered: that \(G_{v+1}\ne0\) for every admissible convex valuation, with its leading parent independent of the valuation, is Corollary~\ref{cor:first-split-anchor} together with the last clause of Proposition~\ref{prop:synchronization}; and that the ratio-identity locus carries at least \(v+1\) positive supported cap shifts through the first-split chain is contained in Corollary~\ref{cor:solid-prefix} below.

\begin{corollary}\label{cor:solid-prefix}
On the ratio-identity locus, under the hypotheses of Theorem~\ref{thm:ratio-rigidity},
\[
G_j\ne0
\qquad(0\le j\le v+1).
\]
\end{corollary}

\begin{proof}
For \(1\le j\le v+1\) this is \eqref{eq:support-prefix}.  The coordinate \(G_0\) is supplied by the virtual row rather than by any genuine one: \(\mathcal M_{-1}^\omega=\{0,1\}\) in \eqref{eq:virtual-edge}, and a column contributes to a minimizer set only where its defect coordinate is nonzero.  The cap-fixed Theorem~\ref{thm:cap-fixed-rigidity} alone does not give \(j=0\), since \(G_0=s_A-s_B\) vanishes on the branch where the two parents are themselves skew-equivalent.
\end{proof}

\subsection{The two floor functionals}\label{subsec:floor-functionals}

\begin{remark}\label{rem:which-floor}
The floor set of Theorem~\ref{thm:floor} is that of \(\alpha_d=\wt(G_d)+d^2\).  A second and equally natural floor functional is the shifted one,
\[
\Phi_r(d)=\wt(G_d)+(r+d)^2
=
\alpha_d+2rd+r^2,
\]
and the two do not have the same minimizers: the term \(2rd\) is not constant in \(d\).  The statement \eqref{eq:floor} does imply the corresponding statement for \(\Phi_r\) when \(r>0\), and the argument is a case split rather than one chain.  Lemma~\ref{lem:laplace-equivalence} gives \(\alpha_0\ge m_+\), so exactly one of two cases holds.  If \(\alpha_0=m_+\), then \(\Phi_r(d)-\Phi_r(0)=(\alpha_d-\alpha_0)+2rd>0\) at every positive shift.  If \(\alpha_0>m_+\), the ratio identity puts \(0\) in the minimizing set.  The premise \(J_0=0\) selects the second case by \eqref{eq:J0-equivalence}, so on the locus the first is vacuous.  At \(r=0\) the two functionals coincide, and on the premise of \eqref{eq:floor} the zeroth shift is then not a minimizer, so the shifted statement does not follow there and is not claimed.  Nothing in this paper uses the passage between the two: Theorem~\ref{thm:two-cap} consumes \eqref{eq:floor-input}, which Corollary~\ref{cor:floor-supplies} supplies directly from \eqref{eq:floor-strong}.  The theorem is named for the implication it proves rather than for the shifted floor, which is a different statement.
\end{remark}

\subsection{An abstract Strict-Monge cofactor principle}\label{subsec:abstract}

The Jacobi--Trudi form is one natural realization of a more general mechanism.  We formulate the abstraction without attempting maximal axiomatic generality.

Let \(R\) be a valued integral domain.  Suppose we have a family of entries \(m(A,d)\in R\) whose support is Ferrers and whose entry valuations \(c(A,d)=\nu(m(A,d))\) are strictly Monge: \(c(A,d)+c(A',d')<c(A,d')+c(A',d)\)
whenever \(A<A'\) and \(d<d'\) and all four entries are supported.  Assume increasing matchings are feasible.

Both boards must be sampled from the \emph{same} family \(m\), which is what
makes a row equal to its parameter; without it the conclusion fails.  Fix an ordered column grid \(d_0<d_1<\cdots<d_n\) and two strictly increasing row-parameter sequences
\[
A=(A_0,\ldots,A_{n-1}),
\qquad
A'=(A'_0,\ldots,A'_{n-1}),
\]
drawn from the index set of \(m\), with
\begin{equation}\label{eq:abstract-prefix}
A_i=A'_i\quad(i<v),
\qquad
A_v\ne A'_v,
\qquad
v\ge1 .
\end{equation}
The two boards are \(M_{ij}=m(A_i,d_j)\) and \(M'_{ij}=m(A'_i,d_j)\); write
\(C_j,C'_j\) for their signed maximal cofactors and \(G_j=C_j-C'_j\).  Because
\(v\ge1\), the two sequences share their top parameter \(A_0\), and it is that
common value which normalizes:
\begin{equation}\label{eq:abstract-normalization}
\sigma_j=\nu(C_j)+c(A_0,d_j),
\qquad
\alpha_j=\nu(G_j)+c(A_0,d_j),
\qquad
m_+=\min_{j>0}\alpha_j,
\end{equation}
and \(\sigma'_j\) likewise.

\begin{theorem}\label{thm:abstract-monge}
With \(m\), \(c\), the grid and \eqref{eq:abstract-prefix} as above, suppose
\(\sigma_0=\sigma'_0\).  Then:
\begin{enumerate}[label=(\roman*),leftmargin=2em]
\item the first forced cofactor-valuation split occurs at column \(v+1\);
\item later columns satisfy a no-backflow inequality in every common-prefix row;
\item exact common-prefix kernel relations force their rightmost valuation minimizers to march strictly right, and those minimizers interlace;
\item if \(\alpha_0>m_+\), the minimizer sets are the adjacent pairs
\[
\{1,2\},\{2,3\},\ldots,\{v,v+1\}.
\]
\end{enumerate}
\end{theorem}

\begin{proof}
Every ingredient is a consequence of the Strict-Monge inequality for \(c\)
together with \eqref{eq:abstract-normalization}, and none uses the ambient
combinatorics.  Theorem~\ref{thm:monge-minor} needs only Ferrers feasibility and
strict Monge, through Lemmas~\ref{lem:strict-monge}
and~\ref{lem:feasible-uncrossing}; the unique initial term of a maximal minor is
then the increasing matching, which is Proposition~\ref{prop:cofactor-initial}.
The action increment of Proposition~\ref{prop:action-recurrence} reads
\[
\sigma_{j+1}-\sigma_j
=
c(A_j,d_j)-c(A_j,d_{j+1})+c(A_0,d_{j+1})-c(A_0,d_j),
\]
which vanishes identically at \(A_j=A_0\), and is positive and strictly
increasing in \(A_j\) beyond it by the Monge inequality at \((A_0,A_j)\) and
\((d_j,d_{j+1})\); here the shared \(A_0\) of
\eqref{eq:abstract-normalization} is exactly what the comparison is taken
against, and the actions, normalized defects and minimizer sets are those of
\eqref{eq:abstract-normalization} with \(c\) in place of \(\omega\).  With \(\sigma_0=\sigma'_0\) in place of
Proposition~\ref{prop:initial-energy} --- the one diagram-specific input, and
here a hypothesis --- Proposition~\ref{prop:synchronization} gives (i).
Lemma~\ref{lem:tail-increase} and Corollary~\ref{cor:first-split-anchor} give
(ii) as in Theorem~\ref{thm:no-backflow}, the latter using only that a strictly
smaller valuation cannot be cancelled.
Lemma~\ref{lem:tropical-balancing} needs only that the exact kernel relation
holds and that the initial forms of a tied minimum cannot sum to zero unless the
tie has at least two members, and Lemma~\ref{lem:marching} then gives (iii),
interlacing included.  Clause (iv) follows as in
Theorem~\ref{thm:cap-fixed-rigidity}: the floor hypothesis excludes column \(0\)
at row \(0\), marching and the no-backflow bound squeeze the rightmost
minimizers to \(R_i=i+2\), and interlacing confines each minimizer set to
\(\{i+1,i+2\}\), where the tie forces equality.
\end{proof}

\begin{remark}\label{rem:abstract-family}
The sampling hypothesis cannot be dropped, and the boards alone do not carry it.
Over \(\Z[t]\) with the \(t\)-adic valuation take
\[
M=\begin{pmatrix}1&1&1\\ t^3&t^2&t\end{pmatrix},
\qquad
N=\begin{pmatrix}1&1&1\\ 2t^3&2t^2&2t\end{pmatrix}.
\]
Both are Ferrers supported with entry valuations \(\binom{0\ 0\ 0}{3\ 2\ 1}\),
which are strictly Monge, and every increasing matching is feasible.  They share
their first row and differ in the second.  But \(C(N)=2\,C(M)\), so
\(\nu(C_j(N))=\nu(C_j(M))\) at every column and there is no split anywhere,
which contradicts (i).  What fails is not Ferrers support or the Monge
inequality but the sampling.  The two second rows carry the same valuation
profile \((3,2,1)\), because multiplying by \(2\) does not change the
\(t\)-adic valuation.  One family \(m(A,d)\) cannot supply both: equal entries
would force equal row parameters, while distinct row parameters with identical
valuation profiles make the defining \(2\times2\) valuation difference vanish
and so violate strict Mongeness.  Hence \(M\) and \(N\) are not two
parameter sequences over one \(m\).  A second board obtained from the first by
scaling a row is the general shape of the obstruction, and it is why the
hypothesis is stated on the parameters rather than on the matrices.
\end{remark}

The Jacobi--Trudi specialization is
\[
R=\Z[h_1,h_2,\ldots],
\qquad
m(A,d)=h_{d-A},
\qquad
c(A,d)=\omega(d-A).
\]

\section{The passive tail}\label{app:rank4}

The results of this appendix concern the columns of the capped cofactor matrix that lie beyond the support of the defect.  None is used in the proof of Theorem~\ref{thm:floor} or Theorem~\ref{thm:boundary-reconstruction}.

Throughout, \(A\) and \(B\) are connected skew diagrams on a common cap grid, \(G=C(A)-C(B)\ne0\), and \(e=\max\supp G\).
Call a column \(j\) \emph{passive} if \(d_j>e\) and \emph{active} otherwise.  Assume throughout this appendix that at least one passive column exists, that is \(e<d_n\); a top-row transfer has \(e=d_n\) and is excluded by the first subsection below.  Write
\[
j_0=\min\{j:d_j>e\},
\qquad
P=n+1-j_0\ \ge1
\]
for the first passive column and the number of passive columns.  Since \(d\) is strictly increasing, the passive columns are the last \(P\).  Call a row \(i\) of one parent \emph{starved} if \(M_{ij}=0\) for every active \(j\), that is if \(d_j<A_i\) whenever \(d_j\le e\), and write \(R_{\mathrm{high}}\) for the set of starved rows.  Since \(A\) is strictly increasing, the starved rows are the last \(|R_{\mathrm{high}}|\).

\subsection{Top-row transfers and passive tails}

A top-row transfer changes the cap origin \(b_0\) and shifts every non-cap \(d_j\).  Connectedness makes the larger parent's terminal maximal cofactor nonzero, so the terminal shift is necessarily supported.  Hence a top-row transfer has no passive tail.  On the ratio-identity locus such a transfer is excluded by Theorem~\ref{thm:terminal-gap}.

\subsection{Starvation and both-passive vanishing}

\begin{lemma}\label{lem:starvation}
\(|R_{\mathrm{high}}|\le P-1\).
\end{lemma}

\begin{proof}
Delete any one passive column.  By Lemma~\ref{lem:cofactor-factorization} the resulting maximal minor is \(\pm s_{\lambda/\nu^{(j)}}\), hence nonzero, so some permutation contributes.  A starved row can be matched only to a passive column, and \(P-1\) passive columns remain; Hall's condition therefore requires \(|R_{\mathrm{high}}|\le P-1\).
\end{proof}

The reverse inequality is not a consequence of the same count and is not proved here, so the vanishing statement below carries the capacity it needs as an explicit hypothesis rather than asserting it of the whole class.  The range over which that capacity has been observed is recorded with the verification scripts.

\begin{corollary}\label{cor:both-passive}
Assume \(P\ge2\) and \(|R_{\mathrm{high}}|=P-1\), and delete two passive columns.  Then every square submatrix of what remains whose row set contains \(R_{\mathrm{high}}\) has determinant zero.
\end{corollary}

\begin{proof}
The remaining matrix has \(P-2\) passive columns, and the \(P-1\) rows of \(R_{\mathrm{high}}\) vanish on every other column.  Hall's condition fails for \(R_{\mathrm{high}}\), so every permutation term contains a zero entry.
\end{proof}

A nonzero square minor retaining every row of \(R_{\mathrm{high}}\) must therefore use at least \(P-1\) passive columns.

\subsection{Mixed factorization}

\begin{proposition}\label{prop:block-factorization}
Let \(S\) be the square submatrix of \(M\) on all \(n\) rows and on a column set \(J\) containing exactly \(|R_{\mathrm{high}}|\) passive columns.  Write \(L\) for the block of \(S\) on the rows outside \(R_{\mathrm{high}}\) and the active columns of \(J\), and \(H\) for the block on \(R_{\mathrm{high}}\) and the passive columns of \(J\).  Then
\begin{equation}\label{eq:block-det}
\det S=\det L\cdot\det H.
\end{equation}
\end{proposition}

\begin{proof}
A starved row vanishes on every active column, so the block of \(S\) on \(R_{\mathrm{high}}\) and the active columns of \(J\) is identically zero.  The starved rows are already the last rows and the passive columns already the last columns, so \(S\) is block upper triangular as it stands and no permutation is needed to reach that form; hence \eqref{eq:block-det} carries no sign.  Both blocks are square because \(J\) was chosen with exactly \(|R_{\mathrm{high}}|\) passive columns.
\end{proof}

So for each admissible column selection the whole passive tail enters through the single scalar \(\det H\), whatever its length, and the mixed minor is the product of that scalar with an active-block determinant.  What \eqref{eq:block-det} does not give is a rank condition on the defect as a whole: that would need identities relating the blocks across different column selections and across the two parents, and none is proved here.

\section*{Declarations}

{\footnotesize
\noindent\textbf{Funding.} This research received no external funding.

\smallskip
\noindent\textbf{Competing interests.} The author has no relevant interests to disclose.

\smallskip
\noindent\textbf{Data availability.} No datasets were generated or analyzed.

\smallskip
\noindent\textbf{Computational verification.} The proofs do not depend on numerical computation.  Scripts used for verification and figures are available at \url{https://github.com/johnnyshields/papers}.

\smallskip
\noindent\textbf{Generative AI use.} During the preparation of this work, the author used OpenAI's ChatGPT (GPT-6 Astra, GPT-5.6 Sol) and Anthropic's Claude (Opus 5.5, Opus 5) to assist with mathematical research, manuscript drafting, and the development of auxiliary verification code.  After using these tools, the author reviewed and edited as needed.  The author takes full responsibility for the content of the work.
\par}

\begin{thebibliography}{9}

\bibitem{Barekat2009CoincidencesCorrigendum}
F. Barekat, V. Reiner and S. van Willigenburg,
\emph{Corrigendum to ``Coincidences among skew Schur functions''},
Adv. Math. \textbf{220} (2009), no.~5, 1655--1656,
DOI: \href{https://doi.org/10.1016/j.aim.2008.10.010}{10.1016/\allowbreak j.aim.2008.10.010}.

\bibitem{Billera2006DecomposableCompositions}
L.~J. Billera, H.~Thomas and S.~van Willigenburg,
\emph{Decomposable compositions, symmetric quasisymmetric functions and equality of ribbon Schur functions},
Adv. Math. \textbf{204} (2006), no.~1, 204--240,
DOI: \href{https://doi.org/10.1016/j.aim.2005.05.014}{10.1016/\allowbreak j.aim.2005.05.014}.

\bibitem{Burkard1996MongeProperties}
R.~E. Burkard, B.~Klinz and R.~Rudolf,
\emph{Perspectives of Monge properties in optimization},
Discrete Appl. Math. \textbf{70} (1996), no.~2, 95--161,
DOI: \href{https://doi.org/10.1016/0166-218X(95)00103-X}{10.1016/\allowbreak 0166-218X(95)00103-X}.

\bibitem{GarciaGarcia2020ToeplitzMinors}
D.~Garc\'ia-Garc\'ia and M.~Tierz,
\emph{Toeplitz minors and specializations of skew Schur polynomials},
J. Combin. Theory Ser. A \textbf{172} (2020), 105201,
DOI: \href{https://doi.org/10.1016/j.jcta.2019.105201}{10.1016/\allowbreak j.jcta.2019.105201}.

\bibitem{Gutschwager2009MultiplicityFree}
C.~Gutschwager,
\emph{Equality of multiplicity free skew characters},
J. Algebraic Combin. \textbf{30} (2009), no.~2, 215--232,
DOI: \href{https://doi.org/10.1007/s10801-008-0158-8}{10.1007/\allowbreak s10801-008-0158-8}.

\bibitem{Jin2024ThickenedRibbon}
E.~Y. Jin and S.~X. Li,
\emph{Towards equivalent thickened ribbon Schur functions},
preprint, \url{https://arxiv.org/abs/2403.01843}.

\bibitem{Jin2025NoncommutingEquality}
E.~Y. Jin and S.~van Willigenburg,
\emph{Equality of skew Schur functions in noncommuting variables},
Bull. Lond. Math. Soc. \textbf{57} (2025), no.~5, 1415--1428,
DOI: \href{https://doi.org/10.1112/blms.70037}{10.1112/\allowbreak blms.70037}.

\bibitem{King2008SkewPositivity}
R.~C. King, T.~A. Welsh and S.~J. van Willigenburg,
\emph{Schur positivity of skew Schur function differences and applications to ribbons and Schubert classes},
J. Algebraic Combin. \textbf{28} (2008), no.~1, 139--167,
DOI: \href{https://doi.org/10.1007/s10801-007-0113-0}{10.1007/\allowbreak s10801-007-0113-0}.

\bibitem{Macdonald1995SymmetricFunctions}
I.~G. Macdonald,
\emph{Symmetric Functions and Hall Polynomials}, 2nd ed.,
Oxford Mathematical Monographs, Oxford University Press, 1995,
DOI: \href{https://doi.org/10.1093/oso/9780198534891.001.0001}{10.1093/\allowbreak oso/\allowbreak 9780198534891.001.0001}.

\bibitem{MarshallOlkinArnold2011Majorization}
A.~W. Marshall, I.~Olkin and B.~C. Arnold,
\emph{Inequalities: Theory of Majorization and Its Applications},
2nd ed., Springer Series in Statistics, Springer, New York, 2011,
DOI: \href{https://doi.org/10.1007/978-0-387-68276-1}{10.1007/\allowbreak 978-0-387-68276-1}.

\bibitem{Maximenko2017TrenchFormulas}
E.~A. Maximenko and M.~A. Moctezuma-Salazar,
\emph{Cofactors and eigenvectors of banded Toeplitz matrices: Trench formulas via skew Schur polynomials},
Oper. Matrices \textbf{11} (2017), no.~4, 1149--1169,
DOI: \href{https://doi.org/10.7153/oam-2017-11-79}{10.7153/\allowbreak oam-2017-11-79}.

\bibitem{McNamara2009Classification}
P.~R.~W. McNamara and S.~van Willigenburg,
\emph{Towards a combinatorial classification of skew Schur functions},
Trans. Amer. Math. Soc. \textbf{361} (2009), no.~8, 4437--4470,
DOI: \href{https://doi.org/10.1090/S0002-9947-09-04683-2}{10.1090/\allowbreak S0002-9947-09-04683-2}.

\bibitem{McNamara2014ComparingSkew}
P.~R.~W. McNamara,
\emph{Comparing skew Schur functions: a quasisymmetric perspective},
J. Comb. \textbf{5} (2014), no.~1, 51--85,
DOI: \href{https://doi.org/10.4310/JOC.2014.v5.n1.a3}{10.4310/\allowbreak JOC.2014.v5.n1.a3}.

\bibitem{OlsonHarris2024HopfApplications}
N.~Olson-Harris,
\emph{Some applications of combinatorial Hopf algebras to integro-differential equations and symmetric function identities},
Ph.D.\ thesis, University of Waterloo, 2024.
\url{https://hdl.handle.net/10012/20710}.

\bibitem{Reiner2007Coincidences}
V.~Reiner, K.~M. Shaw and S.~van Willigenburg,
\emph{Coincidences among skew Schur functions},
Adv. Math. \textbf{216} (2007), no.~1, 118--152,
DOI: \href{https://doi.org/10.1016/j.aim.2007.05.006}{10.1016/\allowbreak j.aim.2007.05.006}.

\bibitem{VanWilligenburg2005Equality}
S.~van Willigenburg,
\emph{Equality of Schur and skew Schur functions},
Ann. Comb. \textbf{9} (2005), no.~3, 355--362,
DOI: \href{https://doi.org/10.1007/s00026-005-0263-5}{10.1007/\allowbreak s00026-005-0263-5}.

\bibitem{Yeats2017HopfSchur}
K.~Yeats,
\emph{A Hopf algebraic approach to Schur function identities},
Electron. J. Combin. \textbf{24} (2017), no.~3, \#P3.10,
DOI: \href{https://doi.org/10.37236/5741}{10.37236/\allowbreak 5741}.

\end{thebibliography}

\end{document}